\documentclass[11pt,oneside]{article}
\usepackage[a4paper,margin=27mm,headheight=15pt]{geometry}
\usepackage[T1]{fontenc}
\usepackage{lmodern,microtype,amsmath,amssymb,amsthm,mathtools,bm}
\usepackage{enumitem,longtable,booktabs,array,xcolor,fancyhdr,xurl,hyperref,bookmark}
\hypersetup{colorlinks=true,linkcolor=blue!45!black,urlcolor=blue!45!black,citecolor=blue!45!black,pdftitle={Discrete Stein--Wainger maximal operator: a forward formalization blueprint},pdfauthor={OpenAI}}
\setlength{\parskip}{4pt}
\setlength{\parindent}{0pt}
\setlist{nosep,leftmargin=2em}
\pagestyle{fancy}\fancyhf{}\fancyhead[L]{\small Discrete Stein--Wainger maximal operator}\fancyhead[R]{\small $1<p<\infty$}\fancyfoot[C]{\thepage}
\renewcommand{\sectionmark}[1]{\markboth{\thesection.\ #1}{}}
\newtheorem{theorem}{Theorem}[section]
\newtheorem{lemma}[theorem]{Lemma}
\newtheorem{proposition}[theorem]{Proposition}
\newtheorem{corollary}[theorem]{Corollary}
\theoremstyle{definition}\newtheorem{definition}[theorem]{Definition}
\newtheorem{construction}[theorem]{Construction}
\theoremstyle{remark}\newtheorem{remark}[theorem]{Remark}
\newcommand{\R}{\mathbb R}\newcommand{\Z}{\mathbb Z}\newcommand{\N}{\mathbb N}\newcommand{\Q}{\mathbb Q}\newcommand{\C}{\mathbb C}\newcommand{\T}{\mathbb T}
\newcommand{\ee}[1]{\mathrm e\!\left(#1\right)}
\newcommand{\norm}[1]{\left\lVert #1\right\rVert}\newcommand{\abs}[1]{\left|#1\right|}
\newcommand{\ip}[2]{\left\langle #1,#2\right\rangle}
\newcommand{\ind}{\mathbf 1}\newcommand{\supp}{\operatorname{supp}}\newcommand{\dist}{\operatorname{dist}}
\newcommand{\lcm}{\operatorname{lcm}}\newcommand{\vol}{\operatorname{vol}}
\newcommand{\av}{\mathop{\mathrm{avg}}}\newcommand{\E}{\mathbb E}\newcommand{\Prb}{\mathbb P}
\newcommand{\CZ}{\mathsf{CZ}}\newcommand{\CSW}{\mathcal C}
\newcommand{\Lift}{\mathcal L}\newcommand{\Aop}{\mathcal A}\newcommand{\Dop}{\mathcal D}
\newcommand{\cF}{\mathcal F}\newcommand{\cS}{\mathcal S}\newcommand{\bM}{\mathsf M}
\newcommand{\bp}[1]{\leavevmode\par\nobreak\noindent{\normalfont\small\sffamily Blueprint identifier: \texttt{#1}.}\par\nobreak}
\newcommand{\uses}[1]{\par\noindent\textsf{Dependencies:} #1.\par}
\newcommand{\ml}[1]{{\normalfont\small\nolinkurl{#1}}}
\newcommand{\bb}[1]{\left\{#1\right\}}
\DeclareMathOperator*{\esssup}{ess\,sup}
\DeclareMathOperator{\sgn}{sgn}
\allowdisplaybreaks[2]
\begin{document}
\hypersetup{pageanchor=false}
\begin{titlepage}
\vspace*{14mm}
{\Large\sffamily MATHEMATICAL FORMALIZATION BLUEPRINT\par}
\vspace{12mm}
{\huge\bfseries The discrete Stein--Wainger\par maximal operator\par}
\vspace{7mm}
{\LARGE The combined $\ell^p$ bound, $1<p<\infty$\par}
\vspace{14mm}
\[
 \left\|\sup_{\lambda\in\R}\left|
 \sum_{y\in\Z^n\setminus\{0\}}
 f(x-y)\,\mathrm e\!\left(\lambda
       (y_1^2+\cdots+y_n^2)^d\right)
 \frac{\Omega(y)}{|y|^n}\right|\right\|_{\ell^p_x}
 \leq C_{p,n,d,\Omega}\|f\|_{\ell^p}.
\]
\vspace{10mm}
Here $n,d\geq1$, $\Omega$ is smooth away from the origin, homogeneous of
 degree zero, and has integral zero on the unit sphere. The parameter
 $\lambda$ is a \emph{scalar}. Only finite exponents $p>1$ are asserted.
\vfill
\textbf{Scope and verification status.}
This document is a mathematical blueprint, not a compiled Lean development.
Its proposed declaration names are local names, not assertions about mathlib's
API. The library boundary is stated in Section~1. Beyond that boundary the
arguments are supplied in the order in which they are used. No harmonic-analysis
or exponential-sum theorem from either paper, or from another paper, is taken
as an axiom.

The starting specifications are the papers
\texttt{arXiv:1907.00405v3} and \texttt{arXiv:2107.14616v2}.
The proof uses equivalent kernel-side decompositions to avoid importing
several of their auxiliary results.
\vspace{8mm}
{\small Prepared 27 September 2026. Mathlib documentation was checked on that
 date; no revision-pinned Lean project has been compiled.}
\end{titlepage}
\hypersetup{pageanchor=true}\pagenumbering{roman}
\tableofcontents
\clearpage\pagenumbering{arabic}

\section{Foundation boundary and conventions}\label{sec:foundation}

\subsection{What is imported}
The following are the mathematical library facts used without reproving them.
They are ordinary mathlib-level results, not placeholders for the difficult
parts of the argument. A ``finite-dimensional Hilbert space'' below may always
be implemented as a finite product of copies of $\C$ with its Euclidean norm.

\begin{description}[style=nextline]
\item[F1. Algebra and finite combinatorics.]
Finite sums and products; induction; pigeonhole counting; binomial and
multinomial identities; the fundamental theorem of arithmetic for positive
integers; gcd, lcm, and the Chinese remainder theorem; polynomial coefficient
identities; a nonzero polynomial has at most its degree many roots.
\item[F2. Scalar and finite-dimensional analysis.]
Completeness of $\R$ and $\C$; real powers; geometric series and elementary
limits; compactness of closed bounded subsets of finite-dimensional spaces;
the extreme-value theorem; differentiation rules; the mean-value theorem,
the fundamental theorem of calculus,
and integration by parts. Change of variables for a differentiable injective
map with nonzero Jacobian is allowed.
\item[F3. Measure and integration.]
Counting measure, Lebesgue measure, product measures, normalized Haar measure
on a finite torus; Tonelli, Fubini, monotone convergence, dominated convergence,
Fatou, measurability of parameter integrals, and the elementary properties of
Bochner integration. Regularity of Lebesgue measure and uniqueness of measures
from a generating family of rectangles are included.
\item[F4. Norms and approximation.]
H\"older's and Minkowski's inequalities, completeness of $L^p$ for $1\leq p<\infty$,
and density of Schwartz functions in Euclidean $L^p$ for $1\leq p<\infty$. Smooth bump functions
exist in finite-dimensional real inner-product spaces. Finite-dimensional
inner-product identities and the adjoint of a bounded Hilbert-space operator
are included. Interpolation, maximal-function estimates, and singular-integral
bounds are \emph{not} included in this item.
\item[F5. Fourier analysis at the library boundary.]
Fourier inversion and differentiation on Schwartz functions, the convolution
identity on that class, Euclidean $L^2$ Plancherel, and the orthonormal Fourier
basis of $L^2(\T^n)$. The resulting unitary identification
$\ell^2(\Z^n)\simeq L^2(\T^n)$ is included. Its elementary consequence
$\|m(D)\|_{2\to2}\leq\|m\|_\infty$ is allowed. No multiplier theorem at
$p\ne2$ is included.
\end{description}

For specificity, the checked documentation includes
\ml{Mathlib.Analysis.Fourier.LpSpace}, with
\ml{MeasureTheory.Lp.norm_fourier_eq};
\ml{Mathlib.Analysis.Fourier.AddCircleMulti}, with
\ml{UnitAddTorus.mFourierBasis} and
\ml{UnitAddTorus.hasSum_sq_mFourierCoeff};
\ml{Mathlib.Analysis.Distribution.SchwartzSpace.Basic}, with
\ml{SchwartzMap.denseRange_toLpCLM}; and
\ml{Mathlib.MeasureTheory.Integral.MeanInequalities}.
The documentation links are collected at the end. These are library entry
points, not a claim that a single declaration implements every conversion
between the conventions below and Lean's bundled types.

\subsection{Types, measures, and Fourier signs}
Fix $n,d\in\N$ with $n,d\geq1$ and put $D=2d$. Use
\[
 E=\R^n,\qquad G=\Z^n,\qquad \T^n=(\R/\Z)^n,
 \qquad P(x)=(x_1^2+\cdots+x_n^2)^d.
\]
The inclusion $G\hookrightarrow E$ is coordinatewise. In particular $P(y)\in\Z$
for $y\in G$. Write $|x|$ for the Euclidean norm, $|x|_\infty$ for the maximum
coordinate norm, and $\mathrm e(t)=\exp(2\pi i t)$. The torus has Haar mass one.
All lattice norms use counting measure.

We deliberately use the \emph{positive} sign for the forward transform:
\begin{align*}
 \widehat f(\xi)&=\sum_{x\in G}f(x)\ee{x\cdot\xi},&
 \check m(y)&=\int_{\T^n}m(\xi)\ee{-y\cdot\xi}\,d\xi,\\
 \widehat h(\xi)&=\int_Eh(x)\ee{x\cdot\xi}\,dx,&
 \check a(y)&=\int_Ea(\xi)\ee{-y\cdot\xi}\,d\xi.
\end{align*}
A Euclidean symbol and its periodization are never identified without saying
so. The notation $m(D)f$ means inverse transform of $m\widehat f$ in the
specified group. Convolution is $(k*f)(x)=\int k(y)f(x-y)\,dy$, or the analogous
sum. Define $M_\beta f(x)=\ee{\beta\cdot x}f(x)$. Thus
\begin{equation}\label{eq:modulation}
 [m(\,\cdot-\beta)](D)=M_{-\beta}m(D)M_\beta.
\end{equation}
This sign convention is equivalent to the library convention by reflection.

\subsection{Quantifiers and constants}
An inequality $A\lesssim B$ means $A\leq CB$ for a finite positive constant
independent of every displayed scale, finite index set, and input function.
Dependencies are stated at the lemma or at the beginning of the section.
Constants may depend on $n,D$ and finitely many seminorms of the fixed smooth
profiles. They may depend on $p$ once an $L^p$ conclusion is asserted.

The notation $\delta^{-O_m(1)}$ has the precise meaning
$C_m\delta^{-C_m}$ for some $C_m\geq1$, independent of $N$ and $\delta\in(0,1/2]$.
It is used only in the finite induction on a polynomial's degree. A finite sum,
product, or fixed positive power of such quantities has the same form: multiply
the constants, add the exponents, and enlarge the larger of the two. A lower
bound $\delta^{O_m(1)}$ means the reciprocal of such an upper bound. These
conventions are existential bounds, not unspecified asymptotic hypotheses.

Every estimate for a supremum is first proved for finite families. Every
lattice $\ell^p$ conclusion can first be read with the output sum restricted to
a finite set. An extended norm is not converted into a real-valued $L^p$ norm
until finiteness has been proved. The limiting rules are proved next.

\section{Elementary norm, maximal, and randomization tools}\label{sec:tools}

\begin{lemma}[Finite reductions and extension]\label{lem:finite}
\bp{finite\_reduction}
Let $1\leq p<\infty$. Finite-support functions are dense in $\ell^p(G)$.
If $u_k(x)\to u(x)$ pointwise and $\|u_k\|_p\leq B$, then $\|u\|_p\leq B$.
For a family $T_a$ on $G$, a bound
$\|\max_{a\in F}|T_af|\|_p\leq B\|f\|_p$ for every finite $F$ implies the same
bound for $\sup_a|T_af|$, provided the latter is pointwise defined as an
extended nonnegative number. On $E$ the same assertion holds for a countable
family. A uniformly bounded family of linear operators on a dense subspace
extends uniquely to its completion.
\end{lemma}
\begin{proof}
For $f_R=f\ind_{\{|x|_\infty\leq R\}}$, the $p$th power of
$\|f-f_R\|_p$ is the tail of the convergent nonnegative series
$\sum_x|f(x)|^p$, hence tends to zero. Fatou gives the second assertion.
For the third, fix a finite output set $E_0$ and a number $B_0$ below
$\sum_{x\in E_0}(\sup_a|T_af(x)|)^p$. At every $x\in E_0$ choose one index
approximating the supremum sufficiently closely that the sum exceeds $B_0$.
Only finitely many indices have been chosen. The assumed bound bounds $B_0$.
Take the supremum over $B_0$, then over finite $E_0$. For a countable family on
$E$, increasing finite maxima and monotone convergence give the assertion.
For extension, send a Cauchy approximating sequence to its image sequence;
the uniform Lipschitz estimate makes the latter Cauchy, and also proves
independence of the approximating sequence. Linearity and the norm estimate
pass to the limit.
\end{proof}

\begin{lemma}[Young, dual testing, and subsequences]\label{lem:young}
\bp{young\_duality\_subsequence}
For $1\leq p\leq\infty$ on $E$ or $G$,
$\|k*f\|_p\leq\|k\|_1\|f\|_p$.
For $1<p<\infty$, the $L^p$ norm of $h$ is the supremum of
$|\int h\overline g|$ over $\|g\|_{p'}\leq1$; bounded functions of finite-measure
support suffice after truncation. If $h_k\to h$ in $L^p$, a subsequence
converges almost everywhere.
\end{lemma}
\begin{proof}
For finite sums Young's inequality follows by applying Minkowski to
$\sum_y k(y)f(\cdot-y)$. For integrals approximate the integrand by simple
functions, use Minkowski, and pass by Fatou; truncating $k$ and $f$ removes any
initial bounded-support assumption. Translation preserves the $L^p$ norm.
The case $p=\infty$ follows directly from the integral of $|k|$.

H\"older gives the upper bound in dual testing. On a finite-measure set where
$1/N\leq|h|\leq N$, put
$g=h|h|^{p-2}/\|h\|_{L^p(\text{that set})}^{p-1}$, with value zero off the
set. Its $L^{p'}$ norm is one and the pairing equals the restricted $L^p$ norm.
An increasing exhaustion and monotone convergence prove equality, including
when the norm is initially infinite. For duality arguments starting on
Schwartz functions, the tests just used can be approximated both in $L^{p'}$
and in $L^2$: truncate, approximate by continuous compactly supported functions,
and convolve with a shrinking smooth bump. Uniform continuity controls the
last step, while regularity of Lebesgue measure controls the preceding step.

Choose a subsequence with $\|h_{k_j}-h\|_p^p\leq2^{-2j p}$.
Chebyshev bounds the measure of $|h_{k_j}-h|>2^{-j}$ by $2^{-jp}$.
The measure of the union of these events for $j\geq J$ tends to zero as
$J\to\infty$. Outside their limsup, the subsequence converges pointwise.
\end{proof}

\begin{lemma}[Distributional interpolation]\label{lem:interpolation}
\bp{sublinear\_interpolation}
Suppose $T$ is homogeneous and subadditive in pointwise norm, possibly with a
finite-dimensional normed input or output. If it is weak type $(a,a)$ with
constant $A$ and weak type $(b,b)$ with constant $B$, where
$1\leq a<p<b<\infty$, then
\[
 \|Tf\|_p\leq C_{a,b,p}A^{1-\theta}B^\theta\|f\|_p,
 \qquad \frac1p=\frac{1-\theta}{a}+\frac\theta b.
\]
If the second endpoint is strong $(\infty,\infty)$, then
$\|Tf\|_p\leq C_{a,p}A^{a/p}B^{1-a/p}\|f\|_p$ for $p>a$.
Strong endpoint estimates may be used in place of weak ones.
\end{lemma}
\begin{proof}
Tonelli applied to $u^p=\int_0^u pt^{p-1}dt$ gives
$\|u\|_p^p=p\int_0^\infty t^{p-1}\mu\{u>t\}\,dt$.
Split $f=f_>+f_\leq$ at $|f|=ct$. Subadditivity gives
\[
 \mu\{|Tf|>t\}\leq (2A/t)^a\!\int_{|f|>ct}|f|^a
 +(2B/t)^b\!\int_{|f|\leq ct}|f|^b.
\]
Substitute this into the layer integral and exchange nonnegative integrals.
The first inner integral is
$\int_0^{|f|/c}t^{p-a-1}dt=(|f|/c)^{p-a}/(p-a)$; the second is
$\int_{|f|/c}^\infty t^{p-b-1}dt=(|f|/c)^{p-b}/(b-p)$.
Consequently
\[
 \|Tf\|_p^p\leq
 \left(\frac{p2^aA^ac^{a-p}}{p-a}
       +\frac{p2^bB^bc^{b-p}}{b-p}\right)\|f\|_p^p.
\]
Choose $c=(A^a/B^b)^{1/(b-a)}$. Taking the $p$th root gives the stated
exponents. Zero endpoint constants are handled by a limit. For the infinite
endpoint choose $c=(2B)^{-1}$, so $|Tf_\leq|\leq t/2$, and retain only the
first term. Strong bounds imply weak bounds by Chebyshev. Splitting according
to the input's pointwise norm proves the vector version without any change.
\end{proof}

\begin{lemma}[The elementary maximal function]\label{lem:maximal}
\bp{hardy\_littlewood\_maximal}
For $X=E$ or $G$ define $\bM f(x)$ as the supremum of the averages of $|f|$ on
centered axis-parallel cubes. On $G$ use integer radii including zero.
Then
\[
 \mu\{\bM f>t\}\leq C_n t^{-1}\|f\|_1,
 \qquad \|\bM f\|_p\leq C_{n,p}\|f\|_p\quad(1<p\leq\infty).
\]
For $R\geq1$ set $B_R(x)=R^{-n}(1+|x|/R)^{-n-1}$. Both its Euclidean integral
and its lattice sum are bounded independently of $R$, and
\begin{equation}\label{eq:envelope}
 B_R*|f|\leq C_n\bM f,
 \qquad B_R*_{E}B_L\leq C_n B_{\max(R,L)}\quad(R,L\geq1).
\end{equation}
\end{lemma}
\begin{proof}
Given finitely many cubes on which the average exceeds $t$, repeatedly select
one of largest side length and discard every remaining cube intersecting it.
The selected cubes are disjoint. Each discarded cube is contained in the
threefold dilation of its selected cube. The union of the original cubes has
measure at most $3^n$ times the sum of the selected measures (with an equally
bounded counting-measure factor on $G$). That sum is at most
$t^{-1}\|f\|_1$. On $G$ apply this to arbitrary finite subsets of the level
set. On $E$ first apply it to a finite subcover of a compact subset of the
level set and use inner regularity. Averages are measurable parameter
integrals; rational radii suffice, since changing the radius changes the
integral continuously by dominated convergence and cube boundaries have
measure zero. Thus the level set is measurable. The $L^\infty$ bound has
constant one. Lemma~\ref{lem:interpolation} proves the strong bounds.

The cube $|x|_\infty\leq R$ has at most $(2R+1)^n$ lattice points, and
Euclidean volume $(2R)^n$. Divide space into $|x|\leq R$ and
$2^kR<|x|\leq2^{k+1}R$. The contribution of the $k$th annulus to the integral
or sum of $B_R$ is at most $C_n2^{-k}$. The same decomposition of the
convolution gives the first inequality in \eqref{eq:envelope}.

For the second assume $R\geq L$. If $|x|\leq2R$, bound $B_R$ by $R^{-n}$
and integrate $B_L$. If $|x|>2R$, split the convolution into
$|z|\leq|x|/2$ and its complement. On the first part
$B_R(x-z)\leq C B_R(x)$. On the second
$B_L(z)\leq C L|x|^{-n-1}\leq C R|x|^{-n-1}$, and integrate $B_R(x-z)$.
This is the required bound. Interchanging $R,L$ treats the other order.
\end{proof}

\begin{lemma}[Schur, Hilbert--Schmidt, and linearization]\label{lem:ttstar}
\bp{ttstar\_and\_linearization}
If $H$ is a kernel with row $L^1$ bound $A$ and column $L^1$ bound $B$, then
$\|H\|_{2\to2}\leq\sqrt{AB}$. Its operator norm is also at most
$\|H\|_{L^2(X\times X)}$ when that norm is finite.
For a bounded Hilbert-space operator $T$, $\|T\|^2=\|TT^*\|$.
A uniform bound for all pointwise selections from a finite family of linear
operators bounds the maximum of their absolute values.
\end{lemma}
\begin{proof}
Cauchy--Schwarz gives
$|Hg(x)|^2\leq(\int|H(x,y)|dy)\int|H(x,y)||g(y)|^2dy$.
Integrate in $x$ to obtain $AB\|g\|_2^2$. Applying Cauchy--Schwarz without
weights proves the Hilbert--Schmidt assertion. Dual testing gives
$\|T\|=\|T^*\|$, and
$\|T^*g\|^2=\ip{TT^*g}{g}\leq\|TT^*\|\|g\|^2$ proves one inequality;
submultiplicativity proves the other. For linearization, at each point choose
an index attaining the finite maximum, resolving ties by a fixed ordering.
The sets on which each index is selected are measurable. Apply the uniform
operator estimate to this selection. Infinite lattice families then follow
from Lemma~\ref{lem:finite}; separable Euclidean families follow by countable
approximation. The initially finite selections used below are bounded
operators even before their uniform estimate: the sum of the finitely many
absolute convolution operators supplies a finite preliminary bound.
\end{proof}

\begin{lemma}[Rademacher square functions]\label{lem:khintchine}
\bp{rademacher\_and\_vector\_extension}
For finitely many vectors $v_i$ in a Hilbert space and independent uniform
signs $\epsilon_i\in\{-1,1\}$, for every $0<p<\infty$,
\[
 c_p\Big(\sum_i\|v_i\|^2\Big)^{1/2}
 \leq\big(\E\|\sum_i\epsilon_i v_i\|^p\big)^{1/p}
 \leq C_p\Big(\sum_i\|v_i\|^2\Big)^{1/2}.
\]
A scalar linear operator bounded on $L^p$ with norm $A$ therefore acts on
finite $\ell^2$-valued sequences with norm at most $C_pA$.
\end{lemma}
\begin{proof}
Put $S=\sum_i\|v_i\|^2$. If $S=0$, every vector is zero and both
claims are immediate. Assume $S>0$. Expanding the inner product gives
$\E\|\sum\epsilon_i v_i\|^2=S$. In the $2r$th moment every term with a sign
appearing once vanishes. Partition the $2r$ positions according to their
common indices. There are at most $(2r)^{2r}$ partitions. A surviving block
has size $m\geq2$, and
$\sum_i\|v_i\|^m\leq S^{m/2}$: divide by $S^{m/2}$ and use that the resulting
nonnegative numbers have sum one and each is at most one. Dropping the
requirement that different blocks have different indices bounds each
partition's contribution by $S^r$. This proves the upper even-moment bound,
and H\"older gives every smaller positive moment.

For the lower bound put $X=\|\sum\epsilon_i v_i\|^2$.
The fourth-moment expansion gives $\E X^2\leq C S^2$. Since
$S\leq S/2+(\E X^2)^{1/2}\Prb\{X>S/2\}^{1/2}$, that event has probability
at least a positive absolute constant. It follows that
$\E X^{p/2}\geq c_pS^{p/2}$.
For the operator assertion randomize its finite list of inputs, use its
scalar $L^p$ bound for every sign choice, average, and apply the two
inequalities just proved pointwise before integrating.
\end{proof}

\begin{lemma}[Binary partial sums]\label{lem:binary}
\bp{binary\_maximal\_inequality}
Let $u_1,\ldots,u_{2^L}$ be complex numbers. For the intervals
$I_{l,k}=\{k2^l+1,\ldots,(k+1)2^l\}$, $0\leq l\leq L$, one has
\[
 \max_{0\leq J\leq2^L}\Big|\sum_{j=1}^J u_j\Big|
 \leq\sum_{l=0}^L\Big(\sum_{k=0}^{2^{L-l}-1}
                 \big|\sum_{j\in I_{l,k}}u_j\big|^2\Big)^{1/2}.
\]
\end{lemma}
\begin{proof}
The binary expansion of $J$ partitions $\{1,\ldots,J\}$ into at most one
aligned interval of each length $2^l$. The triangle inequality bounds the
partial sum by the sum of the absolute values of these interval sums. At each
length a selected absolute value is at most the square root of the sum of all
squared absolute values. The resulting right side does not depend on $J$.
\end{proof}

\section{Elementary rational approximation and cancellation}\label{sec:arithmetic}
Throughout this section the degree is bounded in advance. For $t\in\R$ put
$\|t\|_{\T}=\min_{a\in\Z}|t-a|$. Distances between rational points are torus
distances. All assertions about a denominator use its positive reduced value.

\begin{lemma}[The polynomial identities used by both $TT^*$ arguments]\label{lem:radialcoeff}
\bp{radial\_polynomial\_coefficients}
The polynomial $P$ has integer coefficients. For $y,z\in\Z^n$ and a positive
integer $q$, $y\equiv z\pmod q$ coordinatewise implies
$P(y)\equiv P(z)\pmod q$. In the variable $t=z_i$, with all other coordinates
fixed, $P(z+h)$ has leading coefficient one and coefficient $Dh_i$ at
$t^{D-1}$. Consequently $A P(z-X)-B P(z-Y)$ has coefficients $A-B$ and
$-D(AX_i-BY_i)$ at these two powers. Also, off zero,
$\nabla P(z)=D|z|^{D-2}z$.
\end{lemma}
\begin{proof}
Sums, squares and the integer power $d$ preserve integer coefficients and
congruences, proving the first assertions. Write
$C=\sum_{l\ne i}(z_l+h_l)^2$. The binomial theorem gives
\[
 ((t+h_i)^2+C)^d=(t+h_i)^{2d}
       +\sum_{a=0}^{d-1}\binom da(t+h_i)^{2a}C^{d-a}.
\]
Every term of the remaining sum has degree at most $2d-2$ in $t$.
The binomial expansion of the first term has the stated top two coefficients.
Taking the indicated linear combination gives the difference formula.
Differentiating $(\sum_l z_l^2)^d$ with respect to $z_i$ gives
$2d z_i(\sum_l z_l^2)^{d-1}=D|z|^{D-2}z_i$.
When $n=1$ the sum defining $C$ is empty and equals zero, so the same formulas
apply without a separate dimension restriction.
\end{proof}

\begin{lemma}[Divisors, separation, and lattice counts]\label{lem:divisor}
\bp{divisor\_bound\_and\_rational\_separation}
For every $\epsilon>0$ there is $C_\epsilon$ such that the number $\tau(q)$ of
positive divisors of $q$ is at most $C_\epsilon q^\epsilon$.
Distinct rational points with denominators $q,r$ are separated by at least
$(qr)^{-1}$. If $r\mid q$, then in an interval of length $L$ the number of
integers $h$ for which $r\mid D h$ is at most $D L/r+D+1$.
\end{lemma}
\begin{proof}
Write $q=\prod p^{a_p}$, so $\tau(q)=\prod(a_p+1)$. For
$p\geq2^{2/\epsilon}$, $a+1\leq2^a\leq p^{\epsilon a/2}$.
There are finitely many smaller primes, and for each of them
$\sup_{a\geq0}(a+1)p^{-\epsilon a}<\infty$, since an exponential eventually
dominates $a+1$. Multiplying these finitely many constants proves the bound.
A nonintegral number $a/q-b/r$ has distance at least $1/(qr)$ from $\Z$:
after multiplication by $qr$ the distance is a positive integer.
Finally $r\mid Dh$ is equivalent to $r/\gcd(r,D)\mid h$. Consecutive solutions
are separated by that integer. Counting an arithmetic progression in an
interval gives at most $L\gcd(r,D)/r+1$, which is bounded as stated.
\end{proof}

\begin{lemma}[Many approximate multiples]\label{lem:multiples}
\bp{many\_multiples\_inverse}
There are absolute constants $c,C>0$ with the following property. Suppose
$0<\eta\leq1$, $0\leq E\leq c\eta$, $N\geq C\eta^{-2}$, and
$H\subset\{1,\ldots,N\}$ has at least $\eta N$ elements. If
$\|h\alpha\|_{\T}\leq E$ for every $h\in H$, there is an integer
$1\leq r\leq2/\eta$ such that
\[
 \|r\alpha\|_{\T}\leq C E\eta^{-2}N^{-1}.
\]
For $H\subset[-c_0N,c_0N]\cap\Z$ the same conclusion holds with constants
depending on $c_0$, provided $|H|\geq\eta N$ and the corresponding enlarged
smallness and size hypotheses hold.
\end{lemma}
\begin{proof}
List $H$ increasingly. The sum of its consecutive gaps is less than $N$.
As $|H|-1\geq\eta N/2$, one gap $r$ is at most $2/\eta$.
Subtracting the two approximation inequalities gives an integer $a$ with
$\gamma=r\alpha-a$ and $|\gamma|\leq2E$. If $\gamma=0$ the conclusion holds.
Otherwise partition $1,\ldots,N$ by residues $b$ modulo $r$. In one such
progression the numbers $b\alpha+t\gamma$ range over an interval of length
at most $M|\gamma|$, where $M$ is the number of terms. The condition that
such a number be within $E$ of an integer selects at most
\[
 (M|\gamma|+3)(2E/|\gamma|+1)
 \leq4EM+6E/|\gamma|+3
\]
values of $t$. Indeed at most $M|\gamma|+3$ intervals
$[a-E,a+E]$ meet the range, and each contains at most
$2E/|\gamma|+1$ terms; the last inequality uses $|\gamma|\leq2E$.
Summing over the $r$ residue classes gives
$\eta N\leq4EN+6rE/|\gamma|+3r$.
Choose, for example, $c\leq1/16$ and enlarge $C$ so that
$3r\leq\eta N/4$. It follows that
$|\gamma|\leq12rE/(\eta N)\leq24E/(\eta^2N)$.
For a signed interval remove $0$, choose the positive or negative half with
at least half the remaining elements, and replace negative integers by their
opposites. Increasing its containing interval to
$\{1,\ldots,\lceil c_0N\rceil\}$ reduces it to the proved assertion.
\end{proof}

\begin{lemma}[Inverse polynomial exponential sum]\label{lem:inverseweyl}
\bp{polynomial\_sum\_inverse}
For each $m\geq1$ there are $C_m,c_m>0$ such that the following holds.
Let $F(t)=\sum_{k=0}^m a_k t^k$ be a real polynomial, $N\geq1$, and
$0<\delta\leq1/2$. If
\[
 \left|\sum_{t=1}^N\ee{F(t)}\right|\geq\delta N,
\]
then there is an integer $1\leq q\leq C_m\delta^{-C_m}$ such that
\begin{equation}\label{eq:inversecoeff}
 \|q a_k\|_{\T}\leq C_m\delta^{-C_m}N^{-k}
 \qquad(1\leq k\leq m).
\end{equation}
The same assertion, with constants also depending on $c_0$, holds for a sum
on an integer interval contained in $[-c_0N,c_0N]$ of length at least
$\delta N$. It also holds when the summands have weights $w(t)$ satisfying
$\max|w|+\sum|w(t+1)-w(t)|\leq c_0$, and the weighted sum has magnitude at
least $\delta N$.
\end{lemma}
\begin{proof}
We prove the first assertion by induction on $m$. Each use of
$\delta^{-O_m(1)}$ in this proof means the explicit existential convention
in Section~\ref{sec:foundation}; finite products and powers only enlarge its
constant. Thus a single $C_m$ is chosen after the finitely many bounds at
stage $m$ have been established, not before an unbounded iteration.

For $m=1$, the geometric-sum identity gives
$|\sum_{t=1}^N\ee{a_1t}|\leq\min(N,2/|1-\ee{a_1}|)$.
The inequality $|1-\ee{u}|=2|\sin\pi u|\geq4\|u\|_\T$ follows from
concavity of sine on $[0,\pi/2]$. Hence $q=1$ works.
Assume the assertion proved for degree $m-1$. If
$N\leq C\delta^{-C}$ for any fixed constant required below, take $q=1$
and enlarge the final exponent so that the right side of
\eqref{eq:inversecoeff} is at least $1/2$ for every $k\leq m$.
We may therefore impose every finite polynomial-in-$\delta^{-1}$ lower bound
on $N$ that occurs in the following argument.

Set $z_t=\ee{F(t)}$ on $1\leq t\leq N$ and zero otherwise. Expanding a square,
\[
 \delta^2N^2\leq\left|\sum z_t\right|^2
 =N+2\Re\sum_{h=1}^{N-1}\sum_{t=1}^{N-h}z_{t+h}\overline{z_t}.
\]
Consequently $\sum_{h=1}^{N-1}|C(h)|\geq\delta^2N^2/4$ once
$N\geq2\delta^{-2}$, where $C(h)=\sum_{t=1}^{N-h}\ee{F(t+h)-F(t)}$.
The shifts with $|C(h)|<\delta^2N/8$ contribute less than
$\delta^2N^2/8$. Since every $|C(h)|\leq N$, at least
$\delta^2N/8$ shifts have $|C(h)|\geq\delta^2N/8$; their overlap lengths
$N-h$ are also at least $\delta^2N/8$.

For each selected shift the induction hypothesis, applied to its overlap
interval, supplies $q_h\leq\delta^{-O_m(1)}$ and in particular
\[
 \|q_h m h a_m\|_\T
       \leq\delta^{-O_m(1)}N^{-(m-1)}.
\]
Pigeonhole the integer $q_h$. For one value $q_*$ this holds for at least
$\delta^{O_m(1)}N$ shifts. Lemma~\ref{lem:multiples}, with
$\alpha=q_*ma_m$, $E=\delta^{-O_m(1)}N^{-(m-1)}$, and
$\eta=\delta^{O_m(1)}$, applies after imposing the permitted lower bound on
$N$. It gives an integer $q_0\leq\delta^{-O_m(1)}$ and an integer $b$ such that
\[
 a_m=b/q_0+\varepsilon,
 \qquad |\varepsilon|\leq\delta^{-O_m(1)}N^{-m}.
\]
Here $q_0$ may be the unreduced integer $m q_*r$; reduction is unnecessary.

It remains to approximate the other coefficients with a common denominator.
Choose a positive number $\zeta=\delta^{O_m(1)}$ sufficiently small that
$m|\varepsilon|N^{m-1}(2\zeta N)\leq\delta/(100\pi)$ and
$2\zeta\leq\delta/10$. Choose $L$, a positive multiple of $q_0$, between
$\zeta N$ and $2\zeta N$; this is possible after enlarging the lower bound
on $N$. Divide $1,\ldots,N$ into full consecutive blocks of length $L$ and
one remainder of length less than $L$. Discard the remainder; its contribution
has magnitude less than $\delta N/10$. On a full block and a fixed residue
class modulo $q_0$, the factor $\ee{b t^m/q_0}$ is constant, because
$(t+q_0)^m-t^m$ is divisible by $q_0$. The factor $\ee{\varepsilon t^m}$
differs from its value at the start of the block by at most $\delta/25$,
using $|\ee{u}-\ee{v}|\leq2\pi|u-v|$ and the preceding choice of $\zeta$.
Freezing these factors changes the sum over all full blocks by at most
$\delta N/25$.

The triangle inequality therefore implies that one block-residue progression
has a sum of the lower-degree phase $F_{<m}(t)=\sum_{k<m}a_kt^k$ whose
magnitude is at least $(\delta/2)$ times its length $B=L/q_0$.
Write this progression as $t=t_0+q_0u$, $1\leq u\leq B$, with $|t_0|\leq2N$.
Its coefficients are
\[
 c_l=q_0^l\sum_{k=l}^{m-1}\binom{k}{l}a_k t_0^{k-l}.
\]
The induction hypothesis supplies $r\leq\delta^{-O_m(1)}$ with
$\|r c_l\|_\T\leq\delta^{-O_m(1)}B^{-l}$ for $1\leq l<m$.
Since $B\geq\delta^{O_m(1)}N$, these errors are at most
$\delta^{-O_m(1)}N^{-l}$. Substituting $u=(t-t_0)/q_0$ in the polynomial
identity gives, by the binomial theorem,
\[
 a_k=\sum_{l=k}^{m-1}\binom{l}{k}q_0^{-l}(-t_0)^{l-k}c_l
 \quad(1\leq k<m).
\]
Multiply by $q=rq_0^{m-1}$. Every coefficient multiplying $r c_l$ is an
integer. Its magnitude is at most $\delta^{-O_m(1)}N^{l-k}$, so its error
contribution is at most $\delta^{-O_m(1)}N^{-k}$. Summing at most $m$ such
terms proves \eqref{eq:inversecoeff} for $k<m$. The integer $q$ is a multiple
of $q_0$, so the leading-coefficient estimate also holds. This completes
the induction, including the finite small-$N$ alternatives.

For an interval, translate it to $1,\ldots,L$. Its length satisfies
$\delta N\leq L\leq(2c_0+1)N$. The first assertion controls the translated
coefficients at scale $L$. Translating back uses the same triangular binomial
identity, with translation size at most $(c_0+1)N$. Multiplication of each
error of degree $l$ by a number of size $O(N^{l-k})$ gives the asserted
error of degree $k$ at scale $N$.

For weights let $S_u=\sum_{t=a}^u\ee{F(t)}$. The finite identity
\[
 \sum_{t=a}^b w(t)\ee{F(t)}
   =w(b)S_b+\sum_{t=a}^{b-1}(w(t)-w(t+1))S_t
\]
shows that some prefix has $|S_u|\geq\delta N/c_0$ (replace $c_0$ by
$\max(1,c_0)$). Its length is at least its magnitude. Apply the proved
interval assertion with the correspondingly smaller $\delta$.
\end{proof}

\begin{corollary}[Slicing and an exact rational coefficient]\label{lem:weylcor}
\bp{weighted\_slice\_and\_rational\_cancellation}
Let a polynomial phase of degree at most $m$ be summed on a box of side $N$
in $\Z^n$. Suppose its weights have magnitude at most $A$ and variation at
most $A$ on each line in the first coordinate. If the sum has magnitude at
least $\delta A N^n$, one first-coordinate slice satisfies the conclusion of
Lemma~\ref{lem:inverseweyl}, with constants depending on $n,m$.
In particular, let $N/2\leq Q\leq N^2$ and consider a complete polynomial
sum of side $Q$. If, on every such slice, the coefficient of one power
$t^k$, $k\geq2$, is an exact rational number whose denominator lies in
$[N^{1/2}/2,N]$, then its normalized magnitude is at most $C_mN^{-c_m}$.
\end{corollary}
\begin{proof}
The triangle inequality over the other $n-1$ coordinates gives a slice with
magnitude at least $\delta A N$. Divide its weights by $A$ and apply the
weighted assertion. For the second statement choose a small $c>0$ and suppose
a slice average exceeds $N^{-c}$. Its inverse denominator $r$ is at most
$C_mN^{cC_m}$ and satisfies
$\|r b/q\|_\T\leq C_mN^{cC_m}Q^{-k}$.
Choose $cC_m<1/8$. For sufficiently large $N$, $r<q$ and the displayed
upper bound is less than $1/N\leq1/q$. But $b/q$ is reduced, and $r<q$
forces $\|rb/q\|_\T\geq1/q$. This contradiction bounds every slice average.
Average their absolute values. The remaining finite range of $N$ is absorbed
by enlarging $C_m$.
\end{proof}

\begin{lemma}[A complete quadratic sum]\label{lem:gauss}
\bp{quadratic\_complete\_sum}
Let $A,B,Q\in\Z$, $Q\geq1$, and let $v$ be the reduced denominator of $A/Q$.
Then
\[
 G=Q^{-1}\sum_{z\bmod Q}\ee{(Az^2+Bz)/Q}
 \quad\hbox{satisfies}\quad |G|\leq(2/v)^{1/2}.
\]
If the linear coefficient is $-2\alpha h$, with $\alpha=a/q$ reduced and
$q\mid Q$, a nonzero $G$ additionally requires $q\mid2vh$.
\end{lemma}
\begin{proof}
In $|G|^2$ put $t=z-w$ and sum first over $w$. The finite geometric sum is
zero unless $Q\mid2At$. Consequently
\[
 |G|^2=Q^{-1}\sum_{t\bmod Q:\,Q\mid2At}\ee{(At^2+Bt)/Q},
 \qquad |G|^2\leq\gcd(2A,Q)/Q\leq2/v.
\]
Translation $z\mapsto z+v$ does not change the quadratic factor: writing
$A/Q=a_0/v$, its phase increment is $2a_0z+a_0v\in\Z$.
The linear factor is multiplied by $\ee{-2avh/q}$. A nonzero sum can be
unchanged by that multiplication only if $2avh/q\in\Z$.
Coprimality of $a,q$ yields the stated divisibility.
\end{proof}

\begin{lemma}[Averaging a periodic phase against a smooth weight]\label{lem:periodic}
\bp{periodic\_riemann\_sum}
Fix a Schwartz function $\varphi$ on $E$, put
$\varphi_L(z)=L^{-n}\varphi(z/L)$, and let $L\geq Q\geq1$.
If $a:\Z^n\to\C$ is $Q$-periodic in every coordinate and $|a|\leq1$, then,
for $h\in\Z^n$,
\begin{align*}
 &\sum_{z\in\Z^n}a(z)\varphi_L(z)\overline{\varphi_L(z+h)}\\
 &\quad=\left(Q^{-n}\sum_{r\bmod Q}a(r)\right)
       \int_E\varphi_L(z)\overline{\varphi_L(z+h)}\,dz
       +O\big((Q/L)B_L(h)\big).
\end{align*}
The absolute value of the integral is at most $C B_L(h)$.
\end{lemma}
\begin{proof}
Write $W_h(z)=\varphi_L(z)\overline{\varphi_L(z+h)}$.
The inequality
$(1+|h|/L)\leq(1+|z|/L)(1+|z+h|/L)$ and the Schwartz bounds show, for
arbitrarily large fixed $A$, that
\[
 |W_h(z)|\leq C L^{-2n}(1+|h|/L)^{-n-1}
 \big[(1+|z|/L)^{-A}+(1+|z+h|/L)^{-A}\big].
\]
For $|\nabla W_h|$ the same right side has an additional factor $L^{-1}$.
One way to verify these inequalities is to split into $|z|\geq|h|/2$ and
$|z+h|\geq|h|/2$, allocate $n+1$ decay powers to that larger factor, and
retain $A$ powers in the other. Integration proves the assertion about the
integral.

For a fixed residue $r$, partition $E$ into cells
$r+Qk+[0,Q)^n$. The fundamental theorem of calculus along line segments gives
\[
 \left|Q^n\sum_kW_h(r+Qk)-\int W_h\right|
 \leq C_n Q\sum_k Q^n\sup_{r+Qk+[0,Q]^n}|\nabla W_h|.
\]
Since $Q\leq L$, the decay envelopes vary by at most a fixed factor on a
cell. Their cell sums are bounded by their integrals times a fixed constant:
enlarge each cell by one adjacent layer, compare its supremum with the
integral on that enlargement, and note the bounded overlap. The last display
is therefore at most $C(Q/L)B_L(h)$. Divide by $Q^n$, multiply by $a(r)$,
and sum the $Q^n$ residues. The main terms and errors are exactly those in
the assertion. Every series used here is absolutely convergent by the same
decay bounds.
\end{proof}

\section{Kernel construction and the real-variable multiplier estimates}\label{sec:real}
Fix a complex-valued $C^\infty$ function $\Omega$ on $E\setminus\{0\}$,
homogeneous of degree zero, with integral zero on the unit sphere. Put
$K(x)=\Omega(x)|x|^{-n}$ for $x\ne0$, and put $K(0)=0$ when evaluating a
lattice kernel. Smooth kernels supported away from zero retain their usual
smooth extension with value zero there.

\begin{lemma}[Polar integration and a dyadic resolution]\label{lem:dyadic}
\bp{polar\_cancellation\_and\_dyadic\_kernel}
There is a nonnegative smooth radial function $\psi$, supported where
$1/2\leq|x|\leq2$, such that, on writing
$K_j(x)=K(x)\psi(2^{-j}x)$ for $j\in\Z$,
\begin{gather}
 \|\partial^aK_j\|_\infty\leq C_a2^{-j(n+|a|)},\qquad
 \|K_j\|_{L^1(E)}\leq C,\qquad \int_E K_j=0,\label{eq:dyadicbounds}\\
 \sum_{j\geq1}K_j(x)=(1-\eta(x))K(x),\label{eq:dyadicresolution}
\end{gather}
where $\eta$ is smooth, radial, equal to one on $|x|\leq1$, and zero on
$|x|\geq2$. For $j\geq1$, also $\sum_{y\in G}|K_j(y)|\leq C$.
For every real $\mu$,
\begin{equation}\label{eq:radialcancel}
 \int_E \ee{\mu P(x)}K_j(x)\,dx=0,
 \qquad
 \int_E P(x)\ee{\mu P(x)}K_j(x)\,dx=0.
\end{equation}
\end{lemma}
\begin{proof}
Here is the needed polar integration formula with its measure construction.
On $S^{n-1}$ define
$\sigma(A)=n\,\vol\{r\theta:\theta\in A,\ 0<r\leq1\}$.
Disjointness of rays and countable additivity of Lebesgue measure make this
a measure. Dilation of Lebesgue measure gives, for $0<a<b$,
\[
 \vol\{r\theta:\theta\in A,\ a<r\leq b\}
     =\frac{b^n-a^n}{n}\sigma(A).
\]
Rectangles $A\times(a,b]$ generate the Borel sets of
$S^{n-1}\times(0,\infty)$. Uniqueness of measures and the measurable inverse
$x\mapsto(x/|x|,|x|)$ therefore give
$dx=r^{n-1}\,dr\,d\sigma(\theta)$. This is the usual sphere area measure:
on the upper hemisphere the map
$(r,u)\mapsto r(u,\sqrt{1-|u|^2})$ has absolute Jacobian
$r^{n-1}/\sqrt{1-|u|^2}$, found by differentiating its columns and subtracting
their multiples from the radial column. Thus the induced angular density is
$du/\sqrt{1-|u|^2}$; the lower hemisphere has the same calculation, and their
common equator has cone measure zero. When $n=1$ the formula is simply the
two half-lines and counting measure on $\{-1,1\}$.

Choose a nonnegative smooth bump on $(1,2)$ and divide it by its positive
integral. Let $a(r)$ be its integral from $r$ to $2$, extended by one to
$r\leq1$ and by zero to $r\geq2$. Put $\eta(x)=a(|x|)$ and
$\psi(x)=\eta(x)-\eta(2x)$. It is smooth also at zero because both terms are
constant in a neighborhood of zero. Monotonicity of $a$ gives $\psi\geq0$.
The support and telescoping identity follow from these definitions.
Differentiating the identity $K(tx)=t^{-n}K(x)$ successively gives
$(\partial^aK)(tx)=t^{-n-|a|}(\partial^aK)(x)$. Compactness of
$1/2\leq|x|\leq2$ bounds each derivative there, so the product and chain rules
give \eqref{eq:dyadicbounds}. Scaling its support bounds the integral of its
absolute value. Counting lattice points in $|y|\leq2^{j+1}$ proves the lattice
bound for $j\geq1$.

For a radial factor $b(r)$ supported away from zero, polar integration gives
\[
 \int_E K(x)b(|x|)\,dx
 =\left(\int_0^\infty b(r)\frac{dr}{r}\right)
    \left(\int_{S^{n-1}}\Omega(\theta)\,d\sigma(\theta)\right)=0.
\]
Apply this first with the radial profile defining $K_j$, then with its
products with $\ee{\mu r^D}$ and $r^D\ee{\mu r^D}$. This proves all the
cancellation statements.
\end{proof}

\begin{lemma}[An integrable-kernel singular-integral estimate]\label{lem:cz}
\bp{cz\_from\_decomposition}
Suppose $H\in L^1(E)$ is continuously differentiable off zero, satisfies
\[
 |H(x)|\leq A|x|^{-n},\qquad |\nabla H(x)|\leq A|x|^{-n-1},
\]
and convolution by $H$ has $L^2$ norm at most $A$. Then it has weak $(1,1)$
norm at most $C_nA$, and $L^p$ norm at most $C_{n,p}A$ for every $1<p<\infty$.
The constants do not depend on $\|H\|_1$.
\end{lemma}
\begin{proof}
Divide $H$ by $A$, treating $A=0$ separately, so it suffices to work with
$A=1$. Start with a Schwartz function $f$ and a level $t>0$. Among the
axis-parallel dyadic cubes select the maximal cubes $Q$ on which
$|Q|^{-1}\int_Q|f|>t$. A cube cannot have arbitrarily large ancestors with
this property, since $f\in L^1$. Dyadic cubes are either disjoint or nested;
hence the maximal ones are disjoint. Their total measure is at most
$t^{-1}\|f\|_1$. The parent of a selected cube has average at most $t$, giving
$|Q|^{-1}\int_Q|f|\leq2^nt$. Outside their union, $|f|\leq t$: shrinking
dyadic cubes have diameter tending to zero, and continuity of $f$ makes their
averages tend to its value.

Put $f_Q=|Q|^{-1}\int_Qf$, let $g=f$ off the selected cubes and $g=f_Q$ on
$Q$, and let $b_Q=(f-f_Q)\ind_Q$. Then
\[
 \|g\|_\infty\leq2^nt,\quad \|g\|_1\leq\|f\|_1,\quad
 \|g\|_2^2\leq2^nt\|f\|_1,\quad
 \int b_Q=0,\quad\sum_Q\|b_Q\|_1\leq2\|f\|_1.
\]
Let $Q^*$ be the cube with the same center $c_Q$ and side length
$10\sqrt n$ times that of $Q$. Its union has measure at most
$C_nt^{-1}\|f\|_1$. For $x\notin Q^*$, cancellation gives
\[
 (H*b_Q)(x)=\int_Q[H(x-y)-H(x-c_Q)]b_Q(y)\,dy.
\]
The segment between $x-y$ and $x-c_Q$ stays at distance at least
$|x-c_Q|/2$ from zero. The fundamental theorem of calculus and the gradient
bound therefore bound the bracket by
$C_n|y-c_Q||x-c_Q|^{-n-1}$. Integration in $x$ outside $Q^*$ is at most
$C_n$: partition into shells of radii $2^k\ell(Q)$; their contributions are
at most $C_n2^{-k}\ell(Q)^{-1}$, which are multiplied by
$|y-c_Q|\leq\sqrt n\ell(Q)$. Consequently
\[
 \int_{E\setminus Q^*}|H*b_Q|\leq C_n\|b_Q\|_1.
\]
The series of $b_Q$ converges in $L^2$, since their supports are disjoint and
$f-g\in L^2$. The bounded $L^2$ operator sends its partial sums to an
$L^2$-convergent sequence. An almost-everywhere convergent subsequence and
Fatou transfer the last $L^1$ estimate to $H*\sum_Qb_Q$ off the union of
$Q^*$. Chebyshev there gives a level-set measure at most
$C_nt^{-1}\|f\|_1$. For the good part, Chebyshev and the $L^2$ bound give
\[
 \vol\{|H*g|>t/2\}\leq4t^{-2}\|g\|_2^2
       \leq C_nt^{-1}\|f\|_1.
\]
Together with the expanded cubes this proves the weak estimate.

Lemma~\ref{lem:interpolation} gives the strong bounds for $1<p<2$; $p=2$
is the hypothesis. The adjoint convolution kernel is
$H^*(x)=\overline{H(-x)}$, which has the same size, gradient and $L^2$ bounds.
For $p>2$ apply the proved result to $H^*$ at $p'<2$ and use
\[
 |\langle H*f,g\rangle|=|\langle f,H^**g\rangle|
 \leq C_{n,p}\|f\|_p\|g\|_{p'}.
\]
Dual testing, justified in Lemma~\ref{lem:young}, yields the desired norm.
Schwartz density extends all strong bounds. For weak $(1,1)$ approximate
an $L^1$ input by Schwartz functions; since this particular $H$ belongs to
$L^1$, Young gives convergence of their outputs in $L^1$, and an
almost-everywhere subsequence transfers the distributional bound. Rescaling
$H$ restores $A$.
\end{proof}

\begin{lemma}[Uniform signed dyadic multipliers and their limit]\label{lem:signedK}
\bp{signed\_dyadic\_cz}
For any finite scalar sequence $|c_j|\leq1$, $j\in\Z$, convolution by
$\sum_jc_jK_j$ is bounded on $L^p(E)$ by $C_p$, for $1<p<\infty$.
The symbol
\[
 m_K(\xi)=\sum_{j\geq1}\widehat K_j(\xi),\qquad m_K(0)=0,
\]
is bounded and defines an $L^p$ multiplier with the same uniform type of
bound. Its kernel on Schwartz inputs is
$\mathcal K(x)=K(x)(1-\eta(x))$.
\end{lemma}
\begin{proof}
At each nonzero $x$ only a bounded number of the annular supports meet.
The size and gradient bounds in Lemma~\ref{lem:cz} follow at once from
\eqref{eq:dyadicbounds}. For the $L^2$ bound cancellation and
$|\ee{u}-1|\leq2\pi|u|$ give
$|\widehat K_j(\xi)|\leq C2^j|\xi|$. If $\xi\ne0$, integrate by parts once
in the direction $\xi/|\xi|$. The derivative has $L^1$ norm at most $C2^{-j}$,
so the other bound is $C(2^j|\xi|)^{-1}$. Thus
\begin{equation}\label{eq:minFourier}
 |\widehat K_j(\xi)|\leq C\min(2^j|\xi|,(2^j|\xi|)^{-1}).
\end{equation}
The sum of these upper bounds over all $j\in\Z$ is bounded: split at the
last $j$ with $2^j|\xi|\leq1$ and sum the two geometric series. At $\xi=0$
every summand is zero. Plancherel supplies the required $L^2$ bound, and
Lemma~\ref{lem:cz} applies to every finite sum.

For $j\geq1$, \eqref{eq:minFourier} gives pointwise convergence of the symbol
series. Its finite sums are uniformly bounded. Dominated convergence on the
Fourier side gives $L^2$ convergence of their outputs for each $L^2$ input.
For a Schwartz input, an almost-everywhere subsequence and Fatou give every
claimed $L^p$ bound. The kernel sums equal $K(1-\eta)$ by telescoping.
Convolution of that kernel with a Schwartz function is absolutely convergent
at each point: the kernel is bounded on compact sets, and its decay at
infinity is multiplied by a rapidly decreasing translate of the input.
Dominated convergence identifies the kernel and multiplier limits.
Density defines the $L^p$ extension.
\end{proof}

\begin{construction}[A positive band-limited averaging kernel]\label{con:phi}
Choose a nonzero real even smooth bump $b$ supported in
$[-1/32,1/32]^n$ and put
\[
 \phi(x)=\frac{|\check b(x)|^2}{\int_E|\check b|^2},\qquad
 \phi_J(x)=2^{-Jn}\phi(2^{-J}x),\qquad
 \tau_J(\xi)=\widehat\phi(2^J\xi).
\]
Then $\phi$ is even, nonnegative and Schwartz, $\int\phi=1$,
$\int x_i\phi(x)\,dx=0$, and $\supp\widehat\phi\subset[-1/16,1/16]^n$.
Fourier inversion, the product-convolution identity for Schwartz functions,
and Plancherel prove these statements; the first moments vanish because their
integrands are odd. In particular, $|\phi_J|*|f|\leq C\bM f$, also for the
sampled lattice kernel when $J\geq0$, by Lemma~\ref{lem:maximal}.
\end{construction}

\begin{lemma}[Signed smoothing errors for truncation]\label{lem:smootherrors}
\bp{signed\_truncation\_smoothing\_errors}
For $J\geq1$ define
\[
 E_J(x)=\sum_{j\geq J}K_j(x)-(\phi_J*\mathcal K)(x).
\]
This is an integrable smooth function of integral zero. With $R=2^J$,
\begin{equation}\label{eq:EJdecay}
 |\partial^a E_J(x)|\leq
 C_aR^{-n-|a|}(1+|x|/R)^{-n-2-|a|}\qquad(|a|=0,1).
\end{equation}
Every finite sum $\sum_Jc_JE_J$, $|c_J|\leq1$, is an $L^p(E)$ convolution
operator with norm at most $C_p$, independently of the finite set of $J$.
\end{lemma}
\begin{proof}
The following identity converges absolutely in $L^1$:
\begin{equation}\label{eq:EJseries}
 E_J=\sum_{j\geq J}(K_j-\phi_J*K_j)-\sum_{1\leq j<J}\phi_J*K_j.
\end{equation}
For $j\geq J$, the fundamental theorem of calculus under translation gives
$\|K_j-\phi_J*K_j\|_1\leq\int|z|\phi_J(z)\,dz\,\|\nabla K_j\|_1
\leq C2^{J-j}$.
For $j<J$, use $\int K_j=0$ to subtract $\phi_J(x)$ inside the convolution;
its $L^1$ norm is at most
$\|\nabla\phi_J\|_1\int|z||K_j(z)|\,dz\leq C2^{j-J}$.
Both geometric series converge. Each term has integral zero, proving that
$E_J\in L^1$ and $\int E_J=0$. The pointwise identity follows as well: the
series defining $\mathcal K$ is dominated by $|K|(1-\eta)$ and is absolutely
integrable against a translate of the Schwartz function $\phi_J$.

We give the spatial bounds, including their tails, because they will be used
for signed families. For $|x|\leq16R$, the terms with $j\geq J$ and their
convolutions have derivatives of order $a$ bounded by
$C2^{-j(n+|a|)}$. Their sum is at most $CR^{-n-|a|}$. For $j<J$ subtract
$\partial^a\phi_J(x)$ and use its gradient bound; summing
$C2^jR^{-n-|a|-1}$ gives the same result.

Now let $r=|x|>16R$. Split the scales in \eqref{eq:EJseries} into
$2^j<r/4$, $r/4\leq2^j\leq4r$, and $2^j>4r$.
In the first group $K_j(x)=0$. Cancellation gives, for any prescribed integer
$N_0$,
\[
 |\partial^a(\phi_J*K_j)(x)|
 \leq C_{N_0}2^jR^{-n-|a|-1}(1+r/R)^{-N_0}.
\]
The sum of $2^j$ in this group is at most $Cr$. Its contribution is therefore
at most $CR^{-n-|a|}(1+r/R)^{1-N_0}$.
In the middle group there are a bounded number of indices, all at least $J$.
Taylor's formula, written as two successive fundamental-theorem-of-calculus
integrals, and the vanishing first moments of $\phi_J$ give
\[
 |\partial^a(K_j-\phi_J*K_j)(x)|
 \leq C\left(\int|z|^2\phi_J(z)\,dz\right)
                  \|\nabla^2\partial^aK_j\|_\infty
 \leq C R^2r^{-n-|a|-2}.
\]
In the last group $K_j(x)=0$ and $|x-z|\geq2^{j-2}$ on its support.
The Schwartz bound on $\partial^a\phi_J$, together with $\|K_j\|_1\leq C$,
bounds its sum by
$CR^{-n-|a|}(1+r/R)^{-N_0}$.
Choose $N_0\geq n+|a|+3$. These three estimates prove \eqref{eq:EJdecay}.
They also justify differentiating all the locally uniformly convergent series.

The bounds just proved imply
$\int |x||E_J(x)|\,dx\leq CR$ and
$\|\nabla E_J\|_1\leq C/R$. Cancellation and integration by parts, exactly
as the two individual inequalities in \eqref{eq:minFourier}, give
$|\widehat E_J(\xi)|\leq C\min(R|\xi|,(R|\xi|)^{-1})$.
Thus every finite signed sum has uniformly bounded multiplier and $L^2$ norm.
For its size bound at a point of radius $r$, sum separately $R\leq r$ and
$R>r$. In the first sum \eqref{eq:EJdecay} is at most
$CR^2r^{-n-2}$, whose sum is at most $Cr^{-n}$; in the second it is at most
$CR^{-n}$, whose sum is also at most $Cr^{-n}$. For the gradient the two
bounds are respectively $CR^2r^{-n-3}$ and $CR^{-n-1}$, yielding
$Cr^{-n-1}$. The sums are finite or convergent geometric series even when
$r<2$. Each finite signed sum is integrable. Lemma~\ref{lem:cz} now gives
its uniform $L^p$ bound.
\end{proof}

\section{Oscillatory integrals and a continuous annular maximal estimate}\label{sec:oscillation}
This section proves the continuous estimates needed later rather than importing
a continuous Stein--Wainger theorem.

\begin{lemma}[Polynomial sublevel and oscillatory integral]\label{lem:oscint}
\bp{polynomial\_oscillatory\_integral}
Fix $m\geq2$ and a bounded interval $I$. If a real polynomial
$F(t)=\sum_{k=0}^m a_kt^k$ has
$L=\max_{1\leq k\leq m}|a_k|\geq1$, then for a continuously differentiable
amplitude supported in $I$,
\[
 \left|\int_I\ee{F(t)}a(t)\,dt\right|
 \leq C_{m,I}L^{-1/m}\left(\|a\|_\infty+\int_I|a'|\right).
\]
The same estimate can be applied to one coordinate of an integral over a
fixed bounded box in $E$, if one nonconstant coefficient on every such slice
has magnitude at least $L$ and the amplitude and its variation on each slice
are bounded uniformly.
\end{lemma}
\begin{proof}
First let $Q$ be a polynomial of degree at most $m-1$ whose largest coefficient
has magnitude at least $L$. If $E_\tau=\{t\in I:|Q(t)|\leq\tau\}$ has measure
$a>0$, choose $m$ points in this set with pairwise separation at least
$a/(4m)$. Greedy selection works: the intervals of that radius around fewer
than $m$ selected points remove measure less than $a/2$.
Lagrange interpolation at those points writes $Q$ as
\[
 Q(t)=\sum_{i=1}^m Q(t_i)
             \prod_{h\ne i}\frac{t-t_h}{t_i-t_h}.
\]
To verify this formula subtract the right side from $Q$; the difference has
degree at most $m-1$ and $m$ distinct roots, hence is zero.
Every numerator coefficient is bounded by a constant depending only on
$m,I$, and every denominator has magnitude at least $(a/(4m))^{m-1}$.
Thus $L\leq C\tau a^{-(m-1)}$, or
$|E_\tau|\leq C(\tau/L)^{1/(m-1)}$.

Apply this to $Q=F'$, whose largest coefficient is at least $L$. On the
complement of $|F'|\leq\tau$, subdivide at the roots of $F'-\tau$, $F'+\tau$
and $F''$. A nonzero polynomial has at most its degree many roots; if $F''$
is identically zero it creates no subdivision. There are at most $3m+2$
intervals, and on each $F'$ has fixed sign, magnitude greater than $\tau$,
and is monotone. Integration by parts there gives boundary terms bounded
by $2\|a\|_\infty/(2\pi\tau)$ and integral terms bounded by
\[
 (2\pi\tau)^{-1}\int|a'|
   +(2\pi)^{-1}\|a\|_\infty\int |F''|/|F'|^2.
\]
The last integral is at most $2/\tau$, since $1/F'$ is monotone on that
interval. Adding the intervals and the sublevel set bounds the original
integral by
$C(\|a\|_\infty+\int|a'|)
[(\tau/L)^{1/(m-1)}+\tau^{-1}]$.
Set $\tau=L^{1/m}$. Both powers are $L^{-1/m}$.
For a box apply this proved one-dimensional inequality on each line and
integrate its bound over the other coordinates by Fubini.
\end{proof}

\begin{lemma}[Radial Fourier estimates]\label{lem:radialFourier}
\bp{radial\_oscillatory\_fourier\_square}
Define
\[
 \Psi(t,\eta)=\int_E\ee{tP(z)+z\cdot\eta}K(z)\psi(z)\,dz,
 \qquad
 \Phi_{j,\mu}(\xi)=\widehat{\ee{\mu P}K_j}(\xi)
                  =\Psi(\mu2^{Dj},2^j\xi).
\]
For $|t|\leq2$, both $\Psi(t,\eta)$ and $\partial_t\Psi(t,\eta)$ are bounded
by $C_A\min(|\eta|,|\eta|^{-A})$ for every fixed $A>0$.
For $|t|\geq1$ there are fixed $0<c<C$ such that each of these two functions
has the following bounds:
\begin{equation}\label{eq:stationaryregions}
 \begin{cases}
 C_A|\eta|\,|t|^{-A},& |\eta|\leq c|t|,\\
 C|t|^{-n/2},& c|t|\leq|\eta|\leq C|t|,\\
 C_A|\eta|^{-A},& |\eta|\geq C|t|.
 \end{cases}
\end{equation}
Consequently, for $\sigma\in\{-1,1\}$, $u\in[1,2]$, and any finite subset
of integers $j$,
\begin{align}
 \sup_\xi\sum_j|\Phi_{j,\sigma u2^{\ell-Dj}}(\xi)|^2
 &\leq C\begin{cases}2^{-n\ell}&\ell\geq0,\\1&\ell<0,\end{cases}\label{eq:radialsquare}\\
 \sup_\xi\sum_j|\partial_u\Phi_{j,\sigma u2^{\ell-Dj}}(\xi)|^2
 &\leq C\begin{cases}2^{(2-n)\ell}&\ell\geq0,\\2^{2\ell}&\ell<0.\end{cases}\label{eq:radialsquarederiv}
\end{align}
\end{lemma}
\begin{proof}
The change of variables $x=2^jz$, homogeneity of $K$ and $P$, gives the
formula for $\Phi$. Equation~\eqref{eq:radialcancel} gives
$\Psi(t,0)=\partial_t\Psi(t,0)=0$. Differentiating in $t$ merely changes the
fixed amplitude to $2\pi iP(z)K(z)\psi(z)$; its support and smoothness bounds
are unchanged. We prove all subsequent estimates for either amplitude.

We first record the integration-by-parts mechanism. On a compact support
where $|\nabla F|\geq L$ and derivatives of $F$ of the orders used are at most
$C_aL$, set $b=\nabla F/|\nabla F|^2$. The quotient and product rules give
$|\partial^a b|\leq C_a/L$. Since
$b\cdot\nabla\ee{F}=2\pi i\ee{F}$, integration by parts replaces an amplitude
$a$ by $-(2\pi i)^{-1}\operatorname{div}(ba)$. Repeating $N$ times gives an
integral bounded by $C_NL^{-N}$ times finitely many derivatives of $a$.
The derivative estimate for $b$ follows inductively: each differentiated
inverse power of $|\nabla F|^2$ contributes a product of derivatives of $F$
in its numerator and a compensating additional inverse power in its
denominator. The assumed upper and lower bounds leave one factor $L^{-1}$.
All boundary terms vanish because the amplitude is compactly supported.

On $1/2\leq|z|\leq2$, $|\nabla P(z)|=D|z|^{D-1}$ is bounded above and away
from zero. If $|\eta|\leq c|t|$ with sufficiently small $c$, the gradient of
$tP(z)+\eta\cdot z$ has magnitude at least a constant times $|t|$.
The preceding mechanism bounds the integral and its first $\eta$ derivatives
by $C_A|t|^{-A}$. Integrating those first derivatives along the segment from
$0$ to $\eta$, and using the zero value at $\eta=0$, gives the first line
of \eqref{eq:stationaryregions}. If $|\eta|\geq C|t|$ with sufficiently large
$C$, the gradient has magnitude at least $|\eta|/2$ and all its derivatives
are $O(|\eta|)$, giving the third line. For bounded $t$, this also gives
arbitrary decay at large $|\eta|$; the zero value and bounded first derivatives
give the stated $O(|\eta|)$ bound at bounded $|\eta|$.

For the middle region write the phase as
$|t|F(z)$, where $F(z)=\sgn(t)P(z)+v\cdot z$ and $v=\eta/|t|$ ranges over a
fixed compact annulus. The equation $\nabla F(z)=0$ has exactly one solution
$z_0$: its direction is prescribed by $v$, and its radius is
$(|v|/D)^{1/(D-1)}$. Thus $|z_0|$ is bounded above and away from zero.
For positive radii $r,s$ in a fixed compact subinterval of $(0,\infty)$,
\begin{align*}
 |r^{D-1}\theta-s^{D-1}\omega|^2
  &=(r^{D-1}-s^{D-1})^2
                 +r^{D-1}s^{D-1}|\theta-\omega|^2,\\
 |r\theta-s\omega|^2
  &=(r-s)^2+rs|\theta-\omega|^2.
\end{align*}
The mean-value theorem for $r^{D-1}$ compares the two right sides from below.
It follows that $|\nabla F(z)|\geq c_1|z-z_0|$ on the amplitude's support.
On a shell $|z-z_0|\asymp\rho$ the vector
$b=\nabla F/|\nabla F|^2$ satisfies
$|\partial^a b|\leq C_a\rho^{-1-|a|}$ by the quotient rule just used.

Put $T=|t|$ and $\rho_0=T^{-1/2}$. A smooth cutoff to
$|z-z_0|\leq2\rho_0$ has integral bounded by $C\rho_0^n$ in absolute value.
Use differences of concentric smooth cutoffs to partition the remainder
into shells with radii $\rho=2^k\rho_0$ up to a fixed constant. Their
amplitudes have derivatives of order $a$ bounded by $C_a\rho^{-|a|}$.
After $N$ integrations by parts with the above vector field, the amplitude
is bounded by $C_NT^{-N}\rho^{-2N}$: one divergence differentiates either
the amplitude or $b$, contributing at most two additional inverse powers of
$\rho$, and this proves the bound by induction on $N$. The shell volume is
at most $C\rho^n$. Choose an integer $N>n/2$. The shell bounds sum to
\[
 C\sum_{k\geq0}T^{-N}(2^kT^{-1/2})^{n-2N}
 \leq CT^{-n/2}.
\]
Together with the inner ball this proves the middle estimate.

For the square estimates put $\eta_j=2^j\xi$ and $T=2^\ell$. At $\xi=0$
every term is zero. For $\xi\ne0$ these radii form a geometric progression.
When $T\geq1$, the indices $cT\leq|\eta_j|\leq2CT$ number at most a fixed
constant, and each squared value is $O(T^{-n})$. Below this interval, the
first line of \eqref{eq:stationaryregions} gives
$C T^{-2A}\sum |\eta_j|^2\leq CT^{2-2A}$. Above it the third line gives
$C\sum |\eta_j|^{-2A}\leq CT^{-2A}$. Choose $A\geq n/2+1$.
For $T<1$ use the bounded-$t$ estimate and split the geometric progression
at radius one. This proves \eqref{eq:radialsquare}, also for $\partial_t\Psi$.
Finally differentiation in $u$ multiplies that function by $\sigma2^\ell$;
squaring this factor gives \eqref{eq:radialsquarederiv}.
\end{proof}

\begin{lemma}[Continuous annular maximal decay]\label{lem:continuousband}
\bp{continuous\_annular\_maximal\_decay}
For $T\geq1$ define
\[
 \mathfrak H_T f(x)=
 \sup_{j\in\Z,\ T\leq|\mu|2^{Dj}\leq2T}
       |(\ee{\mu P}K_j)*f(x)|.
\]
For every $1<p<\infty$ there is $\sigma_p>0$ such that
$\|\mathfrak H_Tf\|_{L^p(E)}\leq C_pT^{-\sigma_p}\|f\|_p$.
One may take $\sigma_p=(2D)^{-1}\min(1/p,1/p')$.
\end{lemma}
\begin{proof}
The bound $|K_j|\leq C2^{-jn}\ind_{|x|\leq2^{j+1}}$ gives
$\mathfrak H_T f\leq C\bM f$, including scales $2^j<1$ on $E$.
It remains to obtain decay on $L^2$. Linearize a finite family, writing the
selected scale and parameter at $x$ as $R_x=2^{j(x)}$ and $\mu_x$.
The $TT^*$ kernel is
\[
 H(x,y)=\int_E\ee{\mu_xP(x-z)-\mu_yP(y-z)}
                  K_{j(x)}(x-z)\overline{K_{j(y)}(y-z)}\,dz.
\]
Put $R=\max(R_x,R_y)$, $r=\min(R_x,R_y)$, $h=y-x$.
Its support requires $|h|\leq2R+2r\leq4R$, and the elementary bound is
$|H(x,y)|\leq CR^{-n}$, obtained by integrating the smaller annular support.

If $R/r\geq4$, center coordinates at the smaller-scale center and divide
by $r$. The integration is then over a fixed box. The coefficient of a
pure power $z_i^D$ in the normalized phase has magnitude at least
$T-2T(r/R)^D\geq T/2$. The normalized amplitude and its variation along
each coordinate are at most $CR^{-n}$, by \eqref{eq:dyadicbounds} and
$r/R\leq1$. Lemma~\ref{lem:oscint} gives
$|H(x,y)|\leq CR^{-n}T^{-1/D}$.

Suppose $R/r<4$. By interchanging $x,y$ and conjugating when necessary,
assume $R_y=r$. Use the variable $w=(x-z)/r$. The phase is
$\mu_xr^DP(w)-\mu_yr^DP(w+h/r)$.
Its coefficient of $w_i^{D-1}$ is $-D\mu_y r^{D-1}h_i$; this follows by
expanding $(\sum w_i^2)^d$, and no other coordinate enters that coefficient.
If $|h|\geq rT^{-1/2}$, choose $i$ with $|h_i|\geq|h|/\sqrt n$.
That coefficient has magnitude at least $cT^{1/2}$, and the amplitude again
has uniformly bounded support and variation after division by $R^{-n}$.
Lemma~\ref{lem:oscint} gives
$|H(x,y)|\leq CR^{-n}T^{-1/(2D)}$.
For the remaining pairs retain the elementary bound, supported where
$|h|\leq RT^{-1/2}$.

We have proved the pointwise upper bound
\[
 |H(x,y)|\leq C T^{-1/(2D)}R^{-n}\ind_{|x-y|\leq4R}
       +C R^{-n}\ind_{|x-y|\leq RT^{-1/2}}.
\]
In the bilinear integral against $|g(x)g(y)|$, split into $R_x\geq R_y$ and
its complement. In the first part $R=R_x$. Integrating in $y$ and dropping
the scale restriction bounds the first term by
$CT^{-1/(2D)}\int |g|\bM g$, and the second by
$CT^{-n/2}\int|g|\bM g$, because its ball has relative volume $T^{-n/2}$.
In the other part interchange $x,y$ and obtain exactly the same bounds.
H\"older and the already proved $L^2$ maximal-function bound yield
$\|TT^*\|_{2\to2}\leq CT^{-1/(2D)}$. Lemma~\ref{lem:ttstar} gives
$\|\mathfrak H_T\|_{2\to2}\leq CT^{-1/(4D)}$ for finite families.

Interpolate this bound with weak $(1,1)$ domination by $\bM$ for $p<2$,
and with the $L^\infty$ bound for $p>2$. The exponents from
Lemma~\ref{lem:interpolation} are precisely the displayed $\sigma_p$.
For a Schwartz input, convolution is continuous in $\mu$ on each compact
parameter interval by dominated convergence. A countable dense parameter
set on each scale therefore gives the full supremum. Increasing finite
maxima, monotone convergence, and density complete the assertion.
\end{proof}

\section{Sampling and finite Fourier bookkeeping}\label{sec:sampling}
\begin{lemma}[Periodization, modulation, and products]\label{lem:fourierbook}
\bp{periodization\_and\_spectral\_orthogonality}
For an integrable compactly supported Euclidean symbol $m$, let
$m_{\mathrm{per}}(\xi)=\sum_{k\in\Z^n}m(\xi-k)$. Its $y$th inverse Fourier
coefficient is $\check m(y)$ for $y\in G$. For smooth compactly supported
$m$ this is a rapidly decreasing lattice sequence. If lattice sequences
$u_i$ are rapidly decreasing and their Fourier supports are $A_i\subset\T^n$,
the Fourier support of their product is contained in $A_1+\cdots+A_l$.
In particular a sum of their product over $G$ is zero when that Minkowski
sum excludes zero. Two such products with disjoint Fourier supports are
orthogonal in $\ell^2$.
\end{lemma}
\begin{proof}
Partition $E$ into translates of a fundamental cube and use
$\ee{-y\cdot k}=1$ to compute the inverse coefficient of $m_{\mathrm{per}}$.
Absolute integrability justifies the partition. Repeated integration by
parts proves rapid decrease when $m$ is smooth. For two rapidly decreasing
sequences, insert Fourier inversion for one factor and interchange an
absolutely convergent sum and integral; this gives
$\widehat{uv}=\widehat u*_{\T^n}\widehat v$. Induction proves the product
statement. Its value at zero is the sum of the product over $G$.
For orthogonality use Parseval and the disjoint Fourier supports.
Complex conjugation reflects a Fourier support through zero; thus the same
argument applies to products containing conjugated factors.
\end{proof}

\begin{lemma}[Sampling, also for a maximal family]\label{lem:sampling}
\bp{euclidean\_to\_lattice\_sampling}
Let $1<p<\infty$. Let $m_u$ be finitely many bounded symbols supported in
$[-1/8,1/8]^n$. Suppose
\[
 \left\|\max_u|m_u(D)F|\right\|_{L^p(E)}\leq A\|F\|_{L^p(E)}.
\]
Their periodizations satisfy the corresponding $\ell^p(G)$ estimate with
constant $C_{n,p}A$. If a single $m$ is supported in
$[-1/(8q),1/(8q)]^n$, then
\[
 \Delta_qm(\xi)=\sum_{\beta\in(q^{-1}\Z/\Z)^n}m_{\mathrm{per}}(\xi-\beta)
\]
has $\ell^p$ multiplier norm at most $C_{n,p}$ times the Euclidean $L^p$
multiplier norm of $m$, independently of $q\geq1$.
\end{lemma}
\begin{proof}
Choose a Schwartz function $v$ whose Fourier transform equals one on
$[-1/8,1/8]^n$ and is supported in $[-1/4,1/4]^n$.
For a finite-support lattice function define
$F(x)=\sum_{z\in G}f(z)v(x-z)$. The sum
$\sup_x\sum_z|v(x-z)|$ is bounded: restrict $x$ to a fundamental cube and
sum its Schwartz tail by shells. Weighted H\"older gives
\[
 |F(x)|^p\leq\left(\sum_z|v(x-z)|\right)^{p-1}
                    \sum_z|v(x-z)||f(z)|^p.
\]
Integrate and use translation invariance to obtain
$\|F\|_{L^p}\leq C_p\|f\|_{\ell^p}$.
Write $H_u=m_u(D)F$. On the support of $m_u$,
$\widehat F(\xi)=\sum_z f(z)\ee{z\cdot\xi}$, so Fourier inversion gives
$H_u(z)=m_{u,\mathrm{per}}(D)f(z)$ at every integer $z$.
These point values are defined by the integrable inverse Fourier integral.

Each $H_u$ has Fourier support in $[-1/8,1/8]^n$ and satisfies $H_u=v*H_u$.
This equality first holds in $L^2$ by Fourier multiplication. Both sides are
continuous: the inverse Fourier integral gives continuity of the first, and
H\"older with translation continuity of the Schwartz function $v$ gives it
for the second. Hence equality holds at every point. Thus
\[
 \max_u|H_u(z)|\leq\int_E|v(z-x)|\max_u|H_u(x)|\,dx.
\]
Weighted H\"older, followed by summation in $z$, bounds this sampling map
$L^p(E)\to\ell^p(G)$ uniformly, using
$\int|v|<\infty$ and $\sup_x\sum_z|v(z-x)|<\infty$.
This proves the first assertion. The same proof is valid by approximation
for bounded measurable symbols; their products with $\widehat F$ are
integrable on the fixed compact support.

The inverse kernel of $\Delta_qm$ is
$q^n\ind_{q\mid y_i\ (1\leq i\leq n)}\check m(y)$, by a finite geometric
sum in every coordinate. Split the input into its $q^n$ residue classes.
On the class $r+qG$, identify $r+qz$ with $z$. The kernel becomes
$q^n\check m(qz)$, the inverse transform of $m(\xi/q)$ sampled on $G$.
That symbol is supported in $[-1/8,1/8]^n$. Its Euclidean $L^p$ multiplier
norm is the same as that of $m$: apply the change of variables $x\mapsto qx$
to input and output, whose norm factors cancel. Apply the first assertion
on each residue class and add their $p$th powers.
\end{proof}

\section{A self-contained rational-frequency multiplier lemma}\label{sec:IW}
The result in this section is the rational-frequency transference estimate
needed by the proof. Its combinatorial and analytic ingredients are proved
here. In particular no Ionescu--Wainger theorem is an input.

\begin{definition}[Pairwise coprime blocks]\label{def:blocks}
Fix integers $k\geq1$, $r\geq2$, and a finite set $S$ of pairwise coprime
integers greater than one. Put $q_* =\max(2,\max S)$ and
$Q_A=\prod_{q\in A}q$ for $A\subset S$, with $Q_\varnothing=1$.
Let $T[q]=(q^{-1}\Z/\Z)^n$. The map
$\prod_{q\in A}T[q]\to T[Q_A]$ that adds its coordinates is a bijection.
Indeed a zero sum, multiplied by $Q_A/q$, forces its $q$th component to be
zero by B\'ezout's identity; injectivity and the equal finite cardinalities
$Q_A^n$ then give bijectivity.

Let $\Sigma_A$ consist of the sums whose $q$th component is nonzero for every
$q\in A$ and is zero outside $A$. Put
\[
 \Sigma_{\leq k}=\bigcup_{|A|\leq k}\Sigma_A,
 \qquad
 \Sigma^{A_0}=\bigcup_{\substack{B\subset S\setminus A_0\\|B|\leq k-|A_0|}}
                         \Sigma_B.
\]
These unions are disjoint. For a finite set $\Lambda\subset\T^n$ write
$T_{m;\Lambda}$ for the lattice multiplier
$\sum_{\beta\in\Lambda}m_{\mathrm{per}}(\xi-\beta)$.
All frequency radii in this section satisfy
\begin{equation}\label{eq:epsseparation}
 0<\varepsilon\leq (100r q_*^{2rk})^{-1}.
\end{equation}
Thus a nonzero sum or difference of at most $2r$ centers from
$\Sigma_{\leq k}$ cannot be canceled by perturbations of size at most
$2\varepsilon$ in each summand. To see this, its denominator divides a product
of at most $2rk$ blocks, and a nonzero coordinate has distance at least
$q_*^{-2rk}$ from an integer, greater than $4r\varepsilon$.
\end{definition}

\begin{lemma}[Superorthogonality from a singleton index]\label{lem:superorth}
\bp{singleton\_superorthogonality}
Let $f_i$ be finitely many Hilbert-space-valued functions. Suppose every
integral
$\int\prod_{j=1}^r\langle f_{i_{2j-1}},f_{i_{2j}}\rangle$ vanishes whenever
one index occurs exactly once. Then
\[
 \left\|\sum_i f_i\right\|_{2r}
 \leq C_r\left\|\left(\sum_i\|f_i\|^2\right)^{1/2}\right\|_{2r}.
\]
If functions $f_A$ are indexed by subsets of a finite set of size at most $k$,
and the analogous vanishing holds whenever one underlying element occurs
in exactly one of $A_1,\ldots,A_{2r}$, then the same inequality holds with
constant $C_{k,r}$ and square sum over $A$.
\end{lemma}
\begin{proof}
Expand the $2r$th moment. Terms whose index partition contains a singleton
are zero. There are at most $(2r)^{2r}$ partitions of the $2r$ positions.
For a surviving partition each block has size $m\geq2$. After taking absolute
values and dropping the requirement that different blocks have different
indices, its pointwise sum is at most a product of factors
$\sum_i\|f_i\|^m$. Each such factor is at most
$(\sum_i\|f_i\|^2)^{m/2}$, by the elementary power inequality proved in
Lemma~\ref{lem:khintchine}. The product is the $r$th power of the square sum.
Integrating and taking a $2r$th root proves the first assertion.

First restrict the subset family to sets of one cardinality $l\leq k$.
Color the underlying finite set independently with $l$ colors. A fixed
$l$-element set is transversal (one element of each color) with probability
$l!/l^l$. Therefore the sum of all $f_A$ is $l^l/l!$ times the average of the
sums over transversal sets. Fix a coloring. Write those functions as an
$l$-coordinate array indexed by one element from each color class. Apply the
first assertion to the first coordinate, treating sums in the other
coordinates as its functions. The required vanishing holds because a
singleton element of that color cannot appear in any other color, even after
expanding those sums. For the second coordinate use vectors indexed by the
first coordinate. Expanding their inner products still leaves a singleton
element of the second color in exactly one factor, so the same vanishing
holds. Continue through all $l$ coordinates. The resulting bound is the
$L^{2r}$ norm of the full square sum, with constant $C_r^l$.
Minkowski for the finite color average proves the claim for this cardinality.
For $l=0$ there is only the empty set and no estimate is needed.
Finally sum the $k+1$ cardinalities; each of their square functions is
pointwise at most the square function of the whole family. The triangle
inequality gives the constant $C_{k,r}$.
\end{proof}

\begin{lemma}[An elementary weighted sunflower inequality]\label{lem:sunflower}
\bp{weighted\_sunflower}
For nonnegative weights $w_A$ on subsets $A$ of a finite set with $|A|\leq k$,
\begin{equation}\label{eq:weightedflower}
 \left(\sum_A w_A\right)^r
 \leq C_{k,r}\sum_{A_0}
 \sum_{\substack{A_i\subset S\setminus A_0\ (1\leq i\leq r)\\
                   A_i\text{ pairwise disjoint},\ |A_0\cup A_i|\leq k}}
                  \prod_{i=1}^r w_{A_0\cup A_i}.
\end{equation}
Empty petals are allowed. Thus a repeated set in a sunflower tuple has that
set as core and all its repeated petals empty.
\end{lemma}
\begin{proof}
A family of distinct $l$-element sets with more than $l!(r-1)^l$ members
contains $r$ distinct sets with pairwise intersections all equal. Prove this
by induction on $l$. The case $l=0$ is vacuous since there is only one empty
set. Take a maximal disjoint subfamily. If it has at least $r$ members, they
are the required sets with empty common intersection. Otherwise its union
has at most $l(r-1)$ elements and meets every set in the family. One element
therefore belongs to more than $(l-1)!(r-1)^{l-1}$ sets. Remove it from those
sets, use the induction hypothesis, and restore it to their common
intersection. This proves the assertion. Consequently any family of
\[
 M=1+\sum_{l=0}^k l!(r-1)^l
\]
distinct sets of size at most $k$ contains such an $r$-tuple.

If $W=\sum w_A=0$ there is nothing to prove. Otherwise choose sets independently
with probabilities $w_A/W$. If an atom has probability at least
$1/(4M^2)$, the event that the first $r$ choices equal that atom has probability
at least $(4M^2)^{-r}$ and is an allowed sunflower tuple. If every atom is
smaller, the probability of a collision among $M$ choices is at most
$\binom M2\sum_A(w_A/W)^2\leq1/8$.
With probability at least $7/8$ those choices are distinct and contain an
$r$-tuple of the kind just proved. A union bound over the at most $M^r$
choices of positions shows that the probability that the first $r$ choices
form a sunflower is at least $(7/8)M^{-r}$. In both cases this probability
is at least $(4M^2)^{-r}$.
A tuple with equal pairwise intersections has a uniquely determined core
$A_0$ and disjoint petals $A_i$; conversely such petals give such a tuple.
Multiply the probability inequality by $W^r$. This is
\eqref{eq:weightedflower} with $C_{k,r}=(4M^2)^r$.
\end{proof}

\begin{lemma}[Two Fourier orthogonality statements]\label{lem:blockorthog}
\bp{block\_denominator\_and\_numerator\_orthogonality}
For each $\alpha\in\Sigma_{\leq k}$ let $f_\alpha$ be rapidly decreasing on
$G$ with Fourier support in $\alpha+[-\varepsilon,\varepsilon]^n$.
Put $f_A=\sum_{\alpha\in\Sigma_A}f_\alpha$ and
$f_{A_0,\alpha}=\sum_{\beta\in\Sigma_{A_0}}f_{\alpha+\beta}$ for
$\alpha\in\Sigma^{A_0}$.
First, the functions $f_A$ satisfy the subset vanishing condition of
Lemma~\ref{lem:superorth}. Second, for a fixed core $A_0$ and disjoint petals
$A_1,\ldots,A_r\subset S\setminus A_0$, the functions
\[
 \prod_{i=1}^r f_{A_0,\alpha_i},\qquad \alpha_i\in\Sigma_{A_i},
\]
have pairwise disjoint Fourier supports as the tuple $(\alpha_1,\ldots,\alpha_r)$
varies, and are therefore pairwise orthogonal in $\ell^2$.
\end{lemma}
\begin{proof}
Expand a mixed product in the first assertion into individual centers. If a
block $q$ occurs in exactly one of its sets, the sum of its signed rational
centers has a nonzero $q$-component: the unique contributing center has a
nonzero component there. The direct-product decomposition in
Definition~\ref{def:blocks} prevents cancellation by other blocks.
The separation inequality \eqref{eq:epsseparation} prevents the small
frequency perturbations from reaching zero. The sum of the product is zero
by Lemma~\ref{lem:fourierbook}.

For the second assertion, a Fourier support of the displayed product lies
in a neighborhood of a sum of its petal centers and $r$ centers supported
on the core. For two different tuples, a component belonging to one petal
has a nonzero difference, since each nonempty petal is disjoint from every
other petal and from the core. It cannot be canceled by the other components.
Empty petals have only the zero choice and cannot distinguish two tuples.
Again \eqref{eq:epsseparation} separates the neighborhoods, now of width
$r\varepsilon$ on each side. Parseval gives the orthogonality.
\end{proof}

\begin{lemma}[A norm description]\label{lem:normdescription}
\bp{rational\_frequency\_norm\_description}
With the notation of Lemma~\ref{lem:blockorthog},
\begin{equation}\label{eq:normdescription}
 \left\|\sum_{\alpha\in\Sigma_{\leq k}}f_\alpha\right\|_{2r}^{2r}
 \leq C_{k,r}\sum_{|A_0|\leq k}
       \left\|(f_{A_0,\alpha})_{\alpha\in\Sigma^{A_0}}
                                      \right\|_{\ell^{2r}(G;\ell^2)}^{2r}.
\end{equation}
\end{lemma}
\begin{proof}
The first assertion of Lemma~\ref{lem:blockorthog} and
Lemma~\ref{lem:superorth} bound the left side by
\[ C\sum_x\left(\sum_A|f_A(x)|^2\right)^r.\]
At each $x$ apply Lemma~\ref{lem:sunflower} with $w_A=|f_A(x)|^2$ and sum in
$x$. Each term becomes
\[\left\|\prod_{i=1}^r f_{A_0\cup A_i}\right\|_2^2.\]
The direct block decomposition gives
$f_{A_0\cup A_i}=\sum_{\alpha_i\in\Sigma_{A_i}}f_{A_0,\alpha_i}$.
Expand this product. By the second orthogonality statement its squared norm
is exactly
\[\sum_{\alpha_i\in\Sigma_{A_i}}\left\|\prod_i f_{A_0,\alpha_i}\right\|_2^2.\]
All resulting quantities are nonnegative. Drop the restriction that the
petals be pairwise disjoint, retaining only $|A_0\cup A_i|\leq k$.
Each $\alpha\in\Sigma^{A_0}$ has a unique petal support, so the enlarged sum
is exactly
$\sum_x(\sum_{\alpha\in\Sigma^{A_0}}|f_{A_0,\alpha}(x)|^2)^r$.
Summing in the core proves \eqref{eq:normdescription}.
\end{proof}

\begin{lemma}[Companion projections]\label{lem:companion}
\bp{companion\_projection\_estimate}
Choose a fixed smooth product cutoff $\theta$, equal to one on $[-1,1]^n$
and zero outside $[-2,2]^n$, and put $\theta_\varepsilon(\xi)=\theta(\xi/\varepsilon)$.
For $2\leq t\leq\infty$,
\begin{equation}\label{eq:companion}
 \left\|(T_{\theta_\varepsilon;\alpha+\Sigma_{A_0}}f)
                  _{A_0,\alpha\in\Sigma^{A_0}}
       \right\|_{\ell^t(G\times\{A_0:|A_0|\leq k\};\ell^2_\alpha)}
 \leq C_{n,k,t}\|f\|_{\ell^t(G)}.
\end{equation}
The outside norm at $t=\infty$ is the supremum in $x,A_0$ of the
$\ell^2_\alpha$ norm.
\end{lemma}
\begin{proof}
At $t=2$, the bumps at different centers of $\Sigma_{\leq k}$ are disjoint.
A center with support $A$ occurs in the coset belonging to $(A_0,\alpha)$
exactly when $A_0\subset A$, with $\alpha$ then uniquely determined.
There are at most $2^k$ such cores. Parseval therefore proves the estimate.

For $t=\infty$ fix a core $A_0$. Inclusion--exclusion on which of its block
components are zero gives
\begin{equation}\label{eq:inclusionexclusion}
 T_{m;\Sigma_{A_0}}=
 \sum_{B\subset A_0}(-1)^{|A_0|-|B|}T_{m;T[Q_B]}.
\end{equation}
It suffices to prove the uniform $\ell^2_\alpha$ estimate for one full-grid
coset $\alpha+T[Q]$, where $Q=Q_B$; the triangle inequality then costs at most
$2^k$. Its kernel is
\[
 h_\alpha(y)=Q^n\ind_{y\in QG}\check\theta_\varepsilon(y)\ee{-\alpha\cdot y}.
\]
The cosets for different $\alpha\in\Sigma^{A_0}$ are disjoint by the direct
block decomposition. Their neighborhoods of radius $2\varepsilon$ are
disjoint by \eqref{eq:epsseparation}. Put
$w(y)=\prod_{i=1}^n(1+\varepsilon^2y_i^2)$.
Multiplication of $h_\alpha$ by this polynomial in $y$ differentiates its
Fourier bumps; it does not enlarge their supports. Thus the sequences
$w h_\alpha$ are pairwise orthogonal in $\ell^2$.

Here are the two numerical bounds used in Cauchy--Schwarz:
\begin{align*}
 \sum_{y\in QG}w(y)^{-2}&\leq C_n(\varepsilon Q)^{-n},\\
 \|w h_\alpha\|_2^2&\leq C_n\varepsilon^nQ^n.
\end{align*}
For the first, factor the sum into one-dimensional sums and compare
$\sum_{l\in\Z}(1+(\varepsilon Ql)^2)^{-2}$ with its integral plus its value
at zero. Since $\varepsilon Q\leq1$, the result is $C/(\varepsilon Q)$.
For the second, $\check\theta_\varepsilon(y)=\varepsilon^n\check\theta(\varepsilon y)$.
The product $w(y)\check\theta(\varepsilon y)$ has arbitrary fixed polynomial
decay in $\varepsilon y$. Its squared sum on $QG$ is at most
$C(\varepsilon Q)^{-n}$ by the same comparison, and the prefactor is
$Q^{2n}\varepsilon^{2n}$. This gives the second bound.

For coefficients $c_\alpha$ of squared sum one, Cauchy--Schwarz on $QG$ and
orthogonality now give
\[
 \sum_y\left|\sum_\alpha c_\alpha h_\alpha(y)\right|
 \leq\left(\sum_{y\in QG}w^{-2}\right)^{1/2}
      \left\|\sum_\alpha c_\alpha w h_\alpha\right\|_2
 \leq C_n.
\]
Pair this kernel with a translate of $f$ and dualize the finite
$\ell^2_\alpha$ norm. This proves the $t=\infty$ bound uniformly in the core.
The map in \eqref{eq:companion} is a linear map with pointwise Hilbert norm,
from measure space $G$ to the product measure space of $x$ and cores.
Lemma~\ref{lem:interpolation} between $2$ and $\infty$ proves every remaining
$t$. It applies with these two different measure spaces because its proof
only uses the input splitting and the output distribution function.
\end{proof}

\begin{theorem}[Rational-frequency transference]\label{thm:IW}
\bp{coprime\_block\_multiplier\_transference}
Assume \eqref{eq:epsseparation}. Let $m$ be a bounded measurable Euclidean
symbol supported in $[-\varepsilon/2,\varepsilon/2]^n$, with $L^{2r}(E)$
multiplier norm at most $A$. For $(2r)'\leq p\leq2r$,
\[
 \|T_{m;\Sigma_{\leq k}}f\|_{\ell^p(G)}
       \leq C_{n,k,r,p} A\|f\|_{\ell^p(G)}.
\]
\end{theorem}
\begin{proof}
First take smooth $m$ supported in $[-\varepsilon,\varepsilon]^n$ and
finite-support $f$. Put $f_\alpha=T_{m;\{\alpha\}}f$.
These are rapidly decreasing and have the Fourier supports required by
Lemma~\ref{lem:normdescription}. For a fixed core and outside component,
\[
 f_{A_0,\alpha}=T_{m;\alpha+\Sigma_{A_0}}f
 =T_{m;\alpha+\Sigma_{A_0}}
               T_{\theta_\varepsilon;\alpha+\Sigma_{A_0}}f.
\]
The factorization follows by multiplying symbols: different center bumps
are disjoint and $\theta_\varepsilon=1$ on the support of $m$.
By \eqref{eq:inclusionexclusion}, the operator $T_{m;\Sigma_{A_0}}$ is a sum
of at most $2^k$ full-grid operators. Their $\ell^{2r}$ norms are at most
$C A$ by Lemma~\ref{lem:sampling}, since
$\varepsilon\leq1/(8Q_B)$ follows from \eqref{eq:epsseparation}.
Thus its norm is at most $C2^kA$.

Modulation gives
$T_{m;\alpha+\Sigma_{A_0}}=M_{-\alpha}T_{m;\Sigma_{A_0}}M_\alpha$.
Apply the vector extension in Lemma~\ref{lem:khintchine} to this one common
operator and the finite list of inputs
$M_\alpha T_{\theta_\varepsilon;\alpha+\Sigma_{A_0}}f$.
Input and output modulations preserve their pointwise square norms.
Then sum the $2r$th powers over cores. Lemma~\ref{lem:companion} and
\eqref{eq:normdescription} give
$\|T_{m;\Sigma_{\leq k}}f\|_{2r}\leq C A\|f\|_{2r}$.

Complex conjugation followed by reflection, $f(x)\mapsto\overline{f(-x)}$,
is an isometry and conjugates the multiplier $m$ to $\overline m$.
Therefore their Euclidean $L^{2r}$ multiplier norms are equal.
The adjoint of $T_{m;\Sigma_{\leq k}}$ has the conjugate symbol and is
$T_{\overline m;\Sigma_{\leq k}}$. Apply the proved $2r$ bound to it and use
dual testing to obtain the bound at $(2r)'$. Interpolation gives the
intermediate exponents.

Finally consider the measurable symbol in the statement. Convolve it in
frequency with a nonnegative smooth approximate identity of integral one,
whose support radius tends to zero and is less than $\varepsilon/2$.
The resulting $m_l$ are smooth and supported in
$[-\varepsilon,\varepsilon]^n$. Their Euclidean multiplier norms are at
most $A$: a frequency translation is conjugation by a modulation, and
Minkowski bounds their average of conjugated operators. The functions $m_l$
converge to $m$ in $L^2(E)$. To justify this without a differentiation
 theorem, approximate the compactly supported $L^2$ function $m$ by a
Schwartz function, use Young on the approximation error, and use uniform
continuity and a tail bound on the Schwartz function.
For finite-support $f$ its torus Fourier transform is bounded, and there
are finitely many rational centers. Thus the corresponding lattice outputs
converge in $\ell^2$ by Parseval. An almost-everywhere subsequence and Fatou
transfer the uniform $\ell^p$ bounds. Finite-support density finishes the
proof.
\end{proof}

\begin{construction}[A good set containing every small denominator]\label{con:goodset}
Fix $\rho>0$. For each integer $N\geq2$, let $N_0=\lceil N^{\rho/4}\rceil$
and put
\[
 Q_0=\prod_{\substack{p\leq N_0\\p\text{ prime}}}
                p^{\lfloor\log_p N\rfloor}.
\]
Omit factors with exponent zero. Take as blocks the integer $Q_0$ when it
is greater than one and, for every prime $N_0<p\leq N$, the integer
$p^{\lfloor\log_p N\rfloor}$. They are pairwise coprime. Put
$k=\lceil4/\rho\rceil+1$ and use these blocks to define
$\mathcal U_N=\Sigma_{\leq k}$.
Every rational point of denominator at most $N$ belongs to $\mathcal U_N$.
Indeed the small-prime part of its denominator divides $Q_0$; its number of
distinct prime factors exceeding $N_0\geq N^{\rho/4}$ is at most $4/\rho$,
by multiplying those factors and comparing with $N$. Its denominator thus
divides a product of at most $k$ blocks. The direct-product decomposition
places its rational point in $\Sigma_A$ for a subset of those blocks.

No estimate for the distribution of primes is required. There are at most
$N_0$ small primes, and each block factor there is at most $N$. Hence
$\log Q_0\leq N_0\log N\leq C_\rho N^{\rho/2}$. Every other block is at most
$N$, so $q_*\leq\exp(C_\rho N^{\rho/2})$. Every center of $\mathcal U_N$ has
denominator at most $q_*^k$. Consequently there is a constant $C_{\rho,r}$
such that
\[
 \varepsilon\leq\exp(-C_{\rho,r}N^{2\rho})
\]
implies \eqref{eq:epsseparation} and all conclusions of
Theorem~\ref{thm:IW}, with constants independent of $N$.
Enlarge $C_{\rho,r}$ to include the finite small range where the displayed
elementary power bounds need adjustment. The construction also works when
$N_0\geq N$, in which case all denominators divide the one small-prime block.
\end{construction}

\section{Minor parameter arcs and the choice of the exponent}\label{sec:minor}
For $j\geq1$ and $\lambda\in\T=\R/\Z$ define the finite-kernel operator
\[
 T_{j,\lambda}f(x)=\sum_{y\in G}\ee{\lambda P(y)}K_j(y)f(x-y).
\]
For a positive integer $M$ put
\[
 H_j=2^{\lfloor M\log_2(j+1)\rfloor},\qquad
 w_j=(j+1)^M2^{-Dj},\qquad
 X_j=\bigcup_{\substack{1\leq q<H_j\\0\leq a<q,\ (a,q)=1}}
       \{\lambda:\|\lambda-a/q\|_\T\leq w_j\}.
\]
For $q=1$ the reduced numerator is $a=0$. A supremum over an empty set is
zero in every nonnegative maximal expression below.

\begin{lemma}[Minor-arc decay with an arbitrarily large logarithmic power]\label{lem:minorL2}
\bp{minor\_parameter\_ttstar}
For every $\kappa>0$ there are integers $M,j_1\geq1$ such that for $j\geq j_1$,
\[
 \left\|\sup_{\lambda\notin X_j}|T_{j,\lambda}f|\right\|_2
             \leq C_\kappa j^{-\kappa}\|f\|_2.
\]
\end{lemma}
\begin{proof}
Choose integers $u>2\kappa+2$ and $v>4\kappa+4$. Put $R=2^j$.
First restrict the input to an integer cube $B$ of side $R$ and linearize a
finite set of parameters outside $X_j$. Its output is supported in a cube
$B^*$ of side at most $5R+2$. Write $\lambda_x$ for the selected parameter at
$x$. The $TT^*$ kernel, for $x,y\in B^*$, is
\[
 H(x,y)=\sum_{z\in B}\ee{\lambda_xP(x-z)-\lambda_yP(y-z)}
                      K_j(x-z)\overline{K_j(y-z)}.
\]
The elementary bound is $|H(x,y)|\leq CR^{-n}$: at most $CR^n$ terms have
magnitude at most $CR^{-2n}$.

Call a pair bad if $|H(x,y)|>R^{-n}j^{-u}$.
Use integer coordinates $z=z_0+t$ on $B$, with $1\leq t_i\leq R$, and write
$X=x-z_0$, $Y=y-z_0$. Every coordinate of $X,Y$ is $O(R)$.
After multiplying the amplitude by $R^{2n}$ it and its first-coordinate
variation are bounded by a fixed constant, using
$\|\partial_1K_j\|_\infty\leq CR^{-n-1}$.
Lemma~\ref{lem:weylcor} and the weighted inverse lemma therefore give, for a
bad pair, an integer $q\leq j^b$ such that
\begin{align*}
 \|q(\lambda_x-\lambda_y)\|_\T&\leq j^bR^{-D},\\
 \|qD(\lambda_xX_1-\lambda_yY_1)\|_\T&\leq j^bR^{-(D-1)}.
\end{align*}
Here $b$ is a fixed positive integer depending on $u,n,D$, not on $j,M,x,y$.
To obtain these two coefficients, expand the phase in $t_1$: its degree-$D$
coefficient is $\lambda_x-\lambda_y$, and its degree-$(D-1)$ coefficient is
$-D(\lambda_xX_1-\lambda_yY_1)$. They do not depend on the other coordinates.
All fixed constants from the inverse estimate are absorbed by increasing
$b$, since $C\leq j^{\lceil\log_2\max(C,1)\rceil}$ for $j\geq2$.
Subtract $D X_1$ times the first approximation from the second. The integer
$D X_1$ multiplies an integer approximant to an integer, so, after increasing
$b$ once more,
\begin{equation}\label{eq:minorh}
 \|qD(x_1-y_1)\lambda_y\|_\T\leq j^bR^{-(D-1)}.
\end{equation}

Fix $y$ and $x'=(x_2,\ldots,x_n)$. Suppose at least $Rj^{-v}$ choices of
$x_1$ give bad pairs. Pigeonhole their choices of $q\leq j^b$.
For one $q$, \eqref{eq:minorh} holds on at least $Rj^{-(v+b)}$ distinct
integers $h=x_1-y_1$ in an interval $[-CR,CR]$.
Lemma~\ref{lem:multiples} applies to $qD\lambda_y$ with
$E=j^bR^{-(D-1)}$ and density a fixed multiple of $j^{-(v+b)}$.
Its hypotheses hold for all sufficiently large $j$, because $2^j$ eventually
exceeds every fixed power of $j$. It produces a positive integer $q'$ and
an integer $a'$ such that
\[
 q'\leq Cj^{v+2b},\qquad
 |\lambda_y-a'/q'|\leq Cj^{3b+2v}R^{-D}
\]
after choosing a real representative and, if necessary, adding an integer.
Here $q'=rqD$; dividing the torus approximation by $q'$ only improves the
stated error. Reduce $a'/q'$ to lowest terms.
Choose an integer $M>3b+2v+10$. For sufficiently large $j$, its reduced
denominator is smaller than $H_j\geq(j+1)^M/2$ and its error is smaller than
$w_j=(j+1)^MR^{-D}$. Thus $\lambda_y\in X_j$, a contradiction.

We have proved that, for each fixed $y,x'$, fewer than $Rj^{-v}$ choices of
$x_1$ are bad. Therefore there are at most $CR^{2n}j^{-v}$ bad pairs.
On all other pairs use $|H|\leq R^{-n}j^{-u}$; on bad pairs use the elementary
bound $CR^{-n}$. Since there are at most $CR^{2n}$ pairs in total,
\[
 \|H\|_{\ell^2(B^*\times B^*)}^2
       \leq C(j^{-v}+j^{-2u}).
\]
Hilbert--Schmidt and $TT^*$ give for the localized linearization a norm at
most $C(j^{-v/4}+j^{-u/2})\leq Cj^{-\kappa}$.

Partition the whole input lattice into disjoint cubes of side $R$.
At any output point, only a number bounded in terms of $n$ of their enlarged
cubes can contribute. Cauchy--Schwarz in those contributions and then
summation in the output give the squared global norm bounded by a fixed
multiple of the sum of the squared localized norms. The input squared norms
add exactly, because the cubes are disjoint. This proves the global bound
for each finite linearization. Lemma~\ref{lem:ttstar} and then
Lemma~\ref{lem:finite} give the maximal assertion. Increase $j_1$ to include
all the size conditions imposed during the proof.
\end{proof}

\begin{construction}[Fixing the final exponent and the major-arc parameters]\label{con:parameters}
From now on fix $1<p<\infty$, let $p'=p/(p-1)$, and choose an integer
$r\geq2$ with $2r>\max(p,p')$. Put
$a_p=\min(2/p,2/p')>0$ and choose $\kappa$ with $\kappa a_p>3$.
Take $M,j_1$ from Lemma~\ref{lem:minorL2}, enlarging $M$ to an integer at
least two if necessary. Enlarging $M$ enlarges $X_j$ for all sufficiently
large $j$, so its complement estimate is preserved after increasing $j_1$.
The pointwise bound
$\sup_\lambda|T_{j,\lambda}f|\leq |K_j|*|f|$ has uniformly bounded strong
$(1,1)$ and $(\infty,\infty)$ constants by \eqref{eq:dyadicbounds}.
Interpolating with the minor-arc $L^2$ bound yields
\begin{equation}\label{eq:minorp}
 \left\|\sup_{\lambda\notin X_j}|T_{j,\lambda}f|\right\|_p
                 \leq C_pj^{-2}\|f\|_p\qquad(j\geq j_1).
\end{equation}
Indeed the decay exponent is $\kappa a_p>3$, which is stronger than two.
Set $\rho=1/(16M)$, and form the good rational sets $\mathcal U_N$ using
Construction~\ref{con:goodset} with this $\rho$ and the already chosen $r$.
All choices have thus been made in the order $p$, $r$, $\kappa$, $M$, $\rho$.
\end{construction}

\section{The rational lifting operators and their arithmetic bound}\label{sec:lifting}
Choose smooth real even product cutoffs $\chi,\widetilde\chi$ with
$0\leq\chi,\widetilde\chi\leq1$, such that
$\chi=1$ on $[-1/2,1/2]^n$, $\supp\chi\subset[-1,1]^n$,
$\widetilde\chi=1$ on $[-1,1]^n$, and
$\supp\widetilde\chi\subset[-2,2]^n$.
For $s\geq1$ put $N_s=2^s$ and choose
\begin{equation}\label{eq:vs}
 v_s=\left\lceil C_*(s+1)(1+2^{s/(2M)})\right\rceil,
 \quad L_s=2^{v_s},\quad
 \chi_s(\xi)=\chi(L_s\xi),\quad
 \widetilde\chi_s(\xi)=\widetilde\chi(L_s\xi).
\end{equation}
Choose $C_*$ sufficiently large that, for every $s\geq1$,
\begin{equation}\label{eq:scalerequirements}
 v_s\geq100(n+D+1)(s+1),\qquad
 4/L_s\leq\exp(-C_{\rho,r}N_s^{2\rho}).
\end{equation}
Such a choice exists: $2\rho=1/(8M)<1/(2M)$, and taking logarithms reduces
the second requirement to an inequality dominated by a fixed multiple of
the defining expression for $v_s$. In particular all bumps centered at
$\mathcal U_{N_s}$ used below are disjoint.
Let
\[
 \mathcal R_s=\{a/q\in\T:(a,q)=1,\ 0\leq a<q,\ N_s/2\leq q<N_s\}.
\]
It has at most $N_s^2$ elements. For $\alpha=a/q\in\mathcal R_s$ and
$\beta=b/q\in T[q]$ define
\begin{align}
 S(\alpha,\beta)&=q^{-n}\sum_{z\bmod q}\ee{\alpha P(z)+\beta\cdot z},\label{eq:Gaussdef}\\
 \Lift_{s,\alpha}[m](\xi)
   &=\sum_{\beta\in T[q]}S(\alpha,\beta)
                        (m\chi_s)_{\mathrm{per}}(\xi-\beta),\label{eq:Liftdef}\\
 \Aop_{s,\alpha}&=\Lift_{s,\alpha}[1],\qquad
 \Dop_s[m](\xi)=\sum_{\beta\in\mathcal U_{N_s}}
                    (m\widetilde\chi_s)_{\mathrm{per}}(\xi-\beta).
\end{align}
Symbols and their operators are distinguished by writing $(D)$ when needed.

\begin{lemma}[The exact lifted kernel]\label{lem:liftkernel}
\bp{lifted\_kernel\_identity}
For every bounded measurable $m$ on the support of $\chi_s$,
\begin{equation}\label{eq:liftkernel}
 \check{\Lift_{s,\alpha}[m]}(y)
             =\ee{\alpha P(y)}\check{(m\chi_s)}(y),\qquad y\in G.
\end{equation}
In particular, if $\varphi_s=\check\chi_s$, the kernel of $\Aop_{s,\alpha}$
is $\ee{\alpha P(y)}\varphi_s(y)$, with
$|\varphi_s(y)|\leq C B_{L_s}(y)$.
\end{lemma}
\begin{proof}
Lemma~\ref{lem:fourierbook} computes the coefficient of each shifted bump.
The remaining finite sum is
\[
 \sum_{\beta\in T[q]}S(\alpha,\beta)\ee{-\beta\cdot y}
 =q^{-n}\sum_{z\bmod q}\ee{\alpha P(z)}
                    \sum_{b\bmod q}\ee{b\cdot(z-y)/q}
 =\ee{\alpha P(y)}.
\]
The inner sum is $q^n$ when $z\equiv y\pmod q$ and zero otherwise, by the
finite geometric-sum formula. Since $P$ is an integer polynomial,
$P(z)\equiv P(y)\pmod q$. This proves the identity, including $y=0$.
Finally $\varphi_s(y)=L_s^{-n}\check\chi(y/L_s)$, so its Schwartz bound gives
the envelope.
\end{proof}

\begin{lemma}[Arithmetic maximal decay]\label{lem:arithmax}
\bp{arithmetic\_maximal\_decay}
There is $\gamma>0$, depending on $n,D$, such that
\begin{equation}\label{eq:arith2}
 \left\|\sup_{\alpha\in\mathcal R_s}|\Aop_{s,\alpha}(D)f|\right\|_2
                   \leq C2^{-\gamma s}\|f\|_2.
\end{equation}
For every $1<t<\infty$ there is $\delta_t>0$ such that for any finite list
of inputs $f_\nu$,
\begin{equation}\label{eq:arithvector}
 \left\|\left(\sum_\nu\sup_{\alpha\in\mathcal R_s}
                 |\Aop_{s,\alpha}(D)f_\nu|^2\right)^{1/2}\right\|_t
 \leq C_t2^{-\delta_t s}
          \left\|(f_\nu)_\nu\right\|_{\ell^t(G;\ell^2)}.
\end{equation}
\end{lemma}
\begin{proof}
Write $N=2^s$, $L=L_s$, and linearize the finite rational family by selecting
$\alpha_x=a_x/q_x$ at $x$. The kernel identity gives its $TT^*$ kernel as
\[
 H(x,y)=\sum_{z\in G}\ee{\alpha_xP(z)-\alpha_yP(z+h)}
                         \varphi_s(z)\overline{\varphi_s(z+h)},
 \qquad h=y-x.
\]
The phase is periodic with period $Q=\lcm(q_x,q_y)$ in every coordinate,
where $N/2\leq Q\leq N^2$. Lemma~\ref{lem:periodic} applies because $L\geq N^2$.
Writing $G_Q$ for its normalized complete phase sum, it gives
\begin{equation}\label{eq:arithH}
 |H(x,y)|\leq C\bigl(|G_Q|+N^2/L\bigr)B_L(h).
\end{equation}

Suppose first that $D\geq4$. On every first-coordinate slice the coefficient
of $z_1^{D-1}$ is $-D\alpha_yh_1$. Outside the set
\[
 E_y=\{h\in G:\gcd(q_y,Dh_1)>N^{1/2}\},
\]
its reduced denominator belongs to $[N^{1/2}/2,N]$.
Corollary~\ref{lem:weylcor} therefore bounds $|G_Q|$ by $CN^{-c}$ for some
$c>0$, uniformly in $\alpha_x$ and the other coordinates.
For any cube of side $R\geq N$, membership in $E_y$ implies
$r\mid Dh_1$ for some divisor $r\mid q_y$ with $r>N^{1/2}$.
Lemma~\ref{lem:divisor}, with exponent $1/8$, gives
\[
 \#(E_y\cap\text{that cube})
 \leq CR^{n-1}\sum_{\substack{r\mid q_y\\r>N^{1/2}}}(DR/r+D+1)
 \leq CR^nN^{-3/8}.
\]

Suppose next that $D=2$. The complete sum factors over coordinates into
quadratic sums of the form in Lemma~\ref{lem:gauss}. Let $v$ be the reduced
denominator of $\alpha_x-\alpha_y$. If $v>N^{1/4}$, the first factor, and
hence $|G_Q|$, is at most $CN^{-1/8}$. If $v\leq N^{1/4}$, a nonzero first
factor requires $q_y\mid2v h_1$. Use the larger exceptional set, independent
of $x$,
\[
 E_y=\bigcup_{1\leq v\leq N^{1/4}}\{h:q_y\mid2v h_1\}.
\]
In a cube of side $R\geq N$ its cardinality is at most
\[
 CR^{n-1}\sum_{v\leq N^{1/4}}(C Rv/N+1)
          \leq CR^nN^{-1/2}.
\]
Outside this set $|G_Q|\leq CN^{-1/8}$ in both cases for $v$.
Thus in either degree case there are $c,b>0$ with
$|G_Q|\leq CN^{-c}+\ind_{E_y}(h)$ and exceptional density at most $CN^{-b}$
in all cubes of side at least $N$.

The lattice sum of $B_L$ is bounded. Its exceptional sum is at most
$CN^{-b}$: partition $h$ into the ball of radius $L$ and shells of radius
$2^kL$, use the just-proved counting bound with $R=C2^kL$, and sum the factors
$2^{-k}$. Moreover $N^2/L\leq N^{-10}$ by \eqref{eq:scalerequirements}.
Equation~\eqref{eq:arithH} consequently has a column sum at most $CN^{-c_0}$,
where $c_0=\min(c,b,10)>0$, uniformly in $y$ and the linearization.
The actual kernel $H$ is Hermitian, being a $TT^*$ kernel, so its row bound
is the same. Schur and $TT^*$ give the operator norm at most $CN^{-c_0/2}$.
This proves \eqref{eq:arith2} with $\gamma=c_0/2$, enlarging the constant for
any finite small range of $N$.

For the vector assertion at $t=2$, square \eqref{eq:arith2} for each input and
sum. At both $t=1$ and $t=\infty$, Minkowski and the kernel identity give
pointwise domination by
$|\varphi_s|*(\sum_\nu|f_\nu|^2)^{1/2}$, whose endpoint norm is uniformly
bounded by Young. Interpolating these bounds with the vector $t=2$ bound
proves \eqref{eq:arithvector}; one possible choice is
$\delta_t=\gamma\min(2/t,2/t')$.
\end{proof}

\begin{lemma}[Factorization and lifted square estimates]\label{lem:liftsquare}
\bp{lifted\_square\_function}
One has the exact symbol factorization
\begin{equation}\label{eq:liftfactor}
 \Lift_{s,\alpha}[m]=\Aop_{s,\alpha}\,\Dop_s[m].
\end{equation}
For a finite family of bounded measurable symbols $m_\nu$,
\begin{equation}\label{eq:liftL2square}
 \left\|\left(\sum_\nu\sup_\alpha
          |\Lift_{s,\alpha}[m_\nu](D)f|^2\right)^{1/2}\right\|_2
 \leq C2^{-\gamma s}
      \left(\sup_\eta\sum_\nu|m_\nu(\eta)|^2\right)^{1/2}\|f\|_2.
\end{equation}
If every signed sum $\sum_\nu\epsilon_\nu m_\nu$, $\epsilon_\nu\in\{-1,1\}$,
has Euclidean $L^{2r}$ multiplier norm at most $A$, then
\begin{equation}\label{eq:liftLpsquare}
 \left\|\left(\sum_\nu\sup_\alpha
          |\Lift_{s,\alpha}[m_\nu](D)f|^2\right)^{1/2}\right\|_p
                  \leq C_p2^{-\delta_p s}A\|f\|_p.
\end{equation}
Here and below the supremum in $\alpha$ is over $\mathcal R_s$.
\end{lemma}
\begin{proof}
Every center in $T[q]$ for $q<N_s$ belongs to $\mathcal U_{N_s}$.
The supports of the enlarged bumps at different centers of that good set
are disjoint, and $\chi_s\widetilde\chi_s=\chi_s$.
Multiplication of the two finite periodized sums therefore leaves exactly
the matching-center terms, giving \eqref{eq:liftfactor}.
Apply \eqref{eq:arithvector} to the inputs $\Dop_s[m_\nu](D)f$.
At $p=2$, Parseval and disjointness of the good-set bumps bound their square
norm by
$(\sup_\eta\sum_\nu|m_\nu(\eta)|^2)^{1/2}\|f\|_2$.
At the prescribed $p$, randomize the list by Lemma~\ref{lem:khintchine}.
Each signed symbol multiplied by $\widetilde\chi_s$ has Euclidean
$L^{2r}$ norm at most $CA$, by Young and the uniform $L^1$ norm of
$\check{\widetilde\chi_s}$. Its support is in
$[-2/L_s,2/L_s]^n$. Theorem~\ref{thm:IW} applies with
$\varepsilon=4/L_s$ by \eqref{eq:scalerequirements} and bounds its good-set
periodization by $C_pA$ on $\ell^p$. Average over the signs and use the
Rademacher lower inequality to recover the square norm. Combining this
with the arithmetic vector bound proves \eqref{eq:liftLpsquare}.
\end{proof}

\begin{lemma}[Positive envelopes for lifted annular kernels]\label{lem:liftenvelope}
\bp{lifted\_annular\_positive\_envelope}
For $j\geq1$, every real $\mu$, and every $\alpha\in\mathcal R_s$,
\[
 |\Lift_{s,\alpha}[\Phi_{j,\mu}](D)f|
 \leq C B_{\max(2^j,L_s)}*|f|\leq C\bM f.
\]
The same pointwise bound holds for the supremum over all those indices at
fixed $s$.
\end{lemma}
\begin{proof}
The Euclidean inverse transform of $\Phi_{j,\mu}\chi_s$ is
$(\ee{\mu P}K_j)*_E\varphi_s$. Its absolute value is at most
$|K_j|*_E|\varphi_s|$, which is bounded by
$C B_{2^j}*_E B_{L_s}\leq C B_{\max(2^j,L_s)}$ by
Lemma~\ref{lem:maximal}. Equation~\eqref{eq:liftkernel} does not change the
absolute value when this kernel is sampled on $G$. Apply the lattice
envelope bound from that lemma. Its constant is independent of every
index in the supremum.
\end{proof}

\section{Lifted maximal truncations}\label{sec:truncation}
\begin{lemma}[The unmodulated lifted maximal truncation]\label{lem:lifttrunc}
\bp{lifted\_maximal\_truncation}
There is $\theta_p>0$ such that
\begin{equation}\label{eq:lifttrunc}
 \left\|\sup_{\alpha\in\mathcal R_s,\ J\geq1}
   \left|\Lift_{s,\alpha}\left[\sum_{1\leq j<J}\widehat K_j\right](D)f\right|
       \right\|_p
       \leq C_p2^{-\theta_p s}\|f\|_p.
\end{equation}
The supremum in $J$ is over integers.
\end{lemma}
\begin{proof}
Let $N=2^s$ and let $Q_s=\lcm(1,\ldots,N-1)$. Its elementary upper bound is
$Q_s\leq\prod_{q<N}q\leq N^N$, so $\log_2Q_s\leq s2^s$.
After all constants in \eqref{eq:vs} are fixed, choose an integer $C_1\geq3$
so large that, with the integer power of two
$J_s=2^{C_1(s+1)}$, for every $s\geq1$,
\begin{equation}\label{eq:longthreshold}
 J_s\geq v_s+10,\quad \sqrt n Q_s\leq2^{J_s-3},\quad
          N^{n+2}Q_s2^{-J_s}\leq2^{-10s}.
\end{equation}
Existence follows by taking logarithms: $J_s$ dominates $v_s$, $s2^s$ and
any fixed multiple of $s$; increasing $C_1$ handles the finite initial range.
This choice is permitted to depend on the already fixed $p,n,D,K$.

For $J\leq J_s$, use the binary inequality of Lemma~\ref{lem:binary}, with
$u_j=\Lift_{s,\alpha}[\widehat K_j](D)f(x)$ and
$L=\log_2J_s=C_1(s+1)$. At each binary level the block symbols are
$m_I=\sum_{j\in I}\widehat K_j$. Every signed sum of those block symbols
has coefficients of magnitude at most one on disjoint dyadic indices.
Lemma~\ref{lem:signedK} gives its uniform Euclidean $L^{2r}$ bound.
Equation~\eqref{eq:liftLpsquare} therefore bounds the square function of
that level, with the supremum in $\alpha$ inside each square, by
$C_p2^{-\delta_p s}\|f\|_p$. The triangle inequality over the $L+1$ levels
bounds the short maximum by $C_p(s+1)2^{-\delta_p s}\|f\|_p$.

For $J>J_s$ write
\[
 \sum_{1\leq j<J}\widehat K_j=m_K-m_{\geq J},\qquad
 m_{\geq J}=\tau_Jm_K+\widehat E_J.
\]
The second identity follows by Fourier transformation of
Lemma~\ref{lem:smootherrors}; it is an identity of bounded symbols almost
everywhere, which is sufficient for the operators.
The whole multiplier $m_K$ is bounded after lifting by the one-element case
of \eqref{eq:liftLpsquare}. The supremum of the lifted $\widehat E_J$ is at
most their square function; every finite signed family has a uniformly
bounded real $L^{2r}$ norm by Lemma~\ref{lem:smootherrors}. The same square
estimate gives $C_p2^{-\delta_p s}\|f\|_p$, first for finite $J$ sets and then
for all of them. It remains to bound the long maximum of
$\Lift_{s,\alpha}[\tau_Jm_K]$.

Let $K_s$ denote the lattice operator with symbol
$(m_K\chi_s)_{\mathrm{per}}$. Its $\ell^p$ norm is at most $C_p$, by
Lemma~\ref{lem:sampling}, Lemma~\ref{lem:signedK}, and Young for the cutoff.
For $\beta\in T[q]$ put $G_\beta=K_sM_\beta f$.
Let $\mathcal V_s=\{0,\ldots,Q_s-1\}^n$ and, for $u\in\mathcal V_s$, define
\[
 I_{\alpha,J,u}=\Lift_{s,\alpha}[\ee{u\cdot\xi}\tau_J(\xi)m_K(\xi)](D)f.
\]
For $J>J_s$ and $|u|\leq\sqrt nQ_s$, the fundamental theorem of calculus
and the Schwartz derivative bounds for $\phi$ give on $G$
\[
 |\phi_J(y)-\phi_J(y-u)|\leq C|u|2^{-J}B_{2^J}(y).
\]
Indeed the segment has length $|u|\leq2^{J-3}$, so its decay envelope is
comparable to the envelope centered at $y$.
All scalar symbols in this paragraph have support in a fundamental cube.
Consequently the product of their periodizations is the periodization of
their product: at a torus point only the matching translate can be nonzero.
Lemma~\ref{lem:fourierbook} then identifies multiplication by $\tau_J$ with
convolution by the sampled $\phi_J$, and multiplication by
$\ee{u\cdot\xi}\tau_J$ with its translate $\phi_J(\cdot-u)$.
It follows, using modulation, that
\begin{equation}\label{eq:Ishift}
 I_{\alpha,J,u}(x)=\sum_{\beta\in T[q]}S(\alpha,\beta)\ee{-\beta\cdot x}
                                    (\phi_J*G_\beta)(x-u).
\end{equation}
All convolutions in this display are on $G$.

Since $|S|\leq1$, there are at most $N^{n+2}$ pairs $(\alpha,\beta)$, and
$\bM$ is bounded on $\ell^p$, the translation bound gives uniformly in $u$
\begin{equation}\label{eq:shifterror}
 \left\|\sup_{\alpha,J>J_s}|I_{\alpha,J,0}-I_{\alpha,J,u}|\right\|_p
 \leq C_pN^{n+2}Q_s2^{-J_s}\|f\|_p
 \leq C_p2^{-10s}\|f\|_p.
\end{equation}
For completeness, apply the bound first to each $G_\beta$, obtaining
$CQ_s2^{-J_s}\bM G_\beta$ pointwise after maximizing $J$; then take its
$\ell^p$ norm, use $\|G_\beta\|_p\leq C_p\|f\|_p$, and sum the indicated
finite number of pairs.

We now average the shifted main expression, without paying for the size of
$Q_s$. All averages below are normalized counting averages on $\mathcal V_s$.
In the $p$th power of its norm, put $z=x-u$ in \eqref{eq:Ishift} and then,
for fixed $z$, put $v=z+u\pmod{Q_s}$. Each denominator $q$ under consideration
divides $Q_s$, so $\ee{-\beta\cdot(z+u)}=\ee{-\beta\cdot v}$ and this is a
bijection of the averaging set. Thus
\begin{align*}
 \av_u\sum_x\sup_{\alpha,J>J_s}|I_{\alpha,J,u}(x)|^p
 &=\av_v\sum_z\sup_{\alpha,J>J_s}|(\phi_J*H_{\alpha,v})(z)|^p,\\
 H_{\alpha,v}(z)&=\sum_{\beta\in T[q]}S(\alpha,\beta)
                                     \ee{-\beta\cdot v}G_\beta(z).
\end{align*}
Nonnegativity of $\phi_J$ implies
$\sup_{\alpha,J}|\phi_J*H_{\alpha,v}|
 \leq C\bM(\sup_\alpha|H_{\alpha,v}|)$.
The maximal-function bound therefore makes the last average at most
$C_p\av_v\sum_z\sup_\alpha|H_{\alpha,v}(z)|^p$.

Reverse the variables by the bijection
$(x,u)\mapsto(z=x-u,v=x\pmod{Q_s})$.
The new inner expression is
\[
 H_{\alpha,v}(x-u)=\sum_{\beta\in T[q]}S(\alpha,\beta)
                          \ee{-\beta\cdot x}G_\beta(x-u)
 =\Lift_{s,\alpha}[\ee{u\cdot\xi}m_K](D)f(x).
\]
The Euclidean multiplier $\ee{u\cdot\xi}m_K$ is an output translate of
$m_K(D)$, so its $L^{2r}$ norm is uniformly bounded in $u$.
The one-element lifted estimate \eqref{eq:liftLpsquare} bounds its maximal
$\ell^p$ norm by $C_p2^{-\delta_p s}\|f\|_p$, independently of $u$.
The averaging argument has consequently proved
\[
 \left(\av_u\|\sup_{\alpha,J>J_s}|I_{\alpha,J,u}|\|_p^p\right)^{1/p}
             \leq C_p2^{-\delta_p s}\|f\|_p.
\]
Use Minkowski on the product of $G$ with the normalized finite averaging
space, and use \eqref{eq:shifterror}, to replace $I_{\alpha,J,u}$ by
$I_{\alpha,J,0}$ in this bound. This proves the long smooth maximum.
All manipulations may first use finite $J$ sets; monotone convergence adds
all $J>J_s$.

Combine the short and long estimates. The elementary bound
$(s+1)2^{-\delta_p s}\leq C_p2^{-\delta_p s/2}$ absorbs the short factor.
For example any $0<\theta_p\leq\min(\delta_p/2,1)$ suffices after adjusting
the constant. This proves \eqref{eq:lifttrunc}.
\end{proof}

\section{Lifted oscillatory bands and their assembly}\label{sec:bands}
For $s\geq1$ and $\ell\in\Z$ define
\[
 \mathcal B_{s,\ell}f=
 \sup_{\substack{\alpha\in\mathcal R_s,\ j\geq1\\
                  2^\ell\leq|\mu|2^{Dj}\leq2^{\ell+1}}}
       |\Lift_{s,\alpha}[\Phi_{j,\mu}](D)f|.
\]
Including both endpoints here only enlarges the bounding operator. In the
actual decomposition we use half-open bands.

\begin{lemma}[A uniform-in-band arithmetic gain]\label{lem:middleband}
\bp{lifted\_band\_uniform\_gain}
There is $b_p>0$ such that for every integer $\ell$,
\[
 \|\mathcal B_{s,\ell}f\|_p\leq C_p2^{-b_p s}\|f\|_p.
\]
\end{lemma}
\begin{proof}
First fix one sign $\sigma$ and finitely many $j$. For $u\in[1,2]$ put
$F_{j,\alpha}(u)=\Lift_{s,\alpha}
 [\Phi_{j,\sigma u2^{\ell-Dj}}](D)f$.
These are continuously differentiable in $u$ for finite-support $f$;
differentiation of their compactly supported inverse Fourier integrals is
justified by the derivative bounds of Lemma~\ref{lem:radialFourier}.
At each lattice point the fundamental theorem of calculus gives
\begin{align*}
 \sup_{j,\alpha,u}|F_{j,\alpha}(u)|^2
 &\leq\sum_j\sup_\alpha|F_{j,\alpha}(1)|^2\\
 &\quad+2\int_1^2
 \left(\sum_j\sup_\alpha|F_{j,\alpha}(u)|^2\right)^{1/2}
 \left(\sum_j\sup_\alpha|F'_{j,\alpha}(u)|^2\right)^{1/2}du.
\end{align*}
To obtain it, apply the identity for the derivative of $|F|^2$ to each
$j,\alpha$, extend the integral to $[1,2]$, bound the selected endpoint by
the sum of all endpoints, and apply Cauchy--Schwarz to the $j$ terms in the
integrand. Summing in $x$ and applying Cauchy--Schwarz once more bounds the
integral by the product of the two square-function $\ell^2$ norms.

Equations~\eqref{eq:liftL2square}, \eqref{eq:radialsquare}, and
\eqref{eq:radialsquarederiv} show that the resulting squared norm is at most
\[
 C2^{-2\gamma s}\|f\|_2^2
 \begin{cases}
 2^{-n\ell}+2^{(1-n)\ell},&\ell\geq0,\\
 1+2^\ell,&\ell<0.
 \end{cases}
\]
Each brace is bounded since $n\geq1$. Taking square roots, adding the two
signs, and passing through finite maxima proves a uniform
$C2^{-\gamma s}$ bound at $p=2$.
Lemma~\ref{lem:liftenvelope} supplies weak $(1,1)$ and strong
$(\infty,\infty)$ bounds independent of $s,\ell$. Interpolation gives the
assertion; $b_p=\gamma\min(2/p,2/p')$ is admissible.
\end{proof}

\begin{lemma}[The high-band analytic gain]\label{lem:highband}
\bp{lifted\_high\_band\_decay}
For $\ell\geq0$,
\[
 \|\mathcal B_{s,\ell}f\|_p
            \leq C_p2^{(n+2)s}2^{-\sigma_p\ell}\|f\|_p,
\]
where $\sigma_p>0$ is the exponent from Lemma~\ref{lem:continuousband}.
\end{lemma}
\begin{proof}
The Euclidean maximal family with symbols $\Phi_{j,\mu}\chi_s$ is obtained
by applying the continuous annular family to $\chi_s(D)F$.
Young bounds this preliminary convolution uniformly. Its maximal $L^p$
norm is therefore at most $C_p2^{-\sigma_p\ell}$ by
Lemma~\ref{lem:continuousband}. All these symbols are supported in
$[-1/8,1/8]^n$, so Lemma~\ref{lem:sampling} transfers this bound to their
periodizations. For a shifted center $\beta$, modulation preserves the input
norm and gives the same maximal bound. In the definition of the lift,
$|S(\alpha,\beta)|\leq1$ and there are at most $2^{(n+2)s}$ pairs
$(\alpha,\beta)$. Bound the rational supremum by the sum over those pairs,
use the triangle inequality for the maximal family and then for its
$\ell^p$ norm. This gives the stated estimate.
\end{proof}

\begin{theorem}[Lifted polynomial partial sums]\label{thm:liftedSW}
\bp{lifted\_polynomial\_maximal\_partial\_sums}
There is $c_p>0$ such that
\begin{equation}\label{eq:liftedSW}
 \left\|\sup_{\substack{\alpha\in\mathcal R_s,\ \mu\in\R\\
                        1\leq a\leq b<\infty}}
 \left|\Lift_{s,\alpha}\left[\sum_{j=a}^b\Phi_{j,\mu}\right](D)f\right|
        \right\|_p
        \leq C_p2^{-c_p s}\|f\|_p.
\end{equation}
The endpoints $a,b$ are integers. Denote the maximal expression on the left
before taking its norm by $\mathcal P_s f$.
\end{theorem}
\begin{proof}
Choose $C_2\geq1$ with $C_2\sigma_p>n+3$ and set
$B_s=\lceil C_2(s+1)\rceil$. Fix $\alpha,\mu,a,b$.
If $\mu\ne0$, assign every $j\in[a,b]$ its unique integer band
\[
 2^\ell\leq |\mu|2^{Dj}<2^{\ell+1}.
\]
At most one $j$ lies in a prescribed band: increasing $j$ by one multiplies
$|\mu|2^{Dj}$ by $2^D\geq4$, whereas two numbers in the same half-open band
have ratio less than two. Split these bands into
$\ell\geq B_s$, $-B_s\leq\ell<B_s$, and $\ell<-B_s$.
The high-band contribution is pointwise at most
$\sum_{\ell\geq B_s}\mathcal B_{s,\ell}f$.
By Lemma~\ref{lem:highband} its norm is at most
\[
 C_p2^{(n+2)s}\sum_{\ell\geq B_s}2^{-\sigma_p\ell}\|f\|_p
 \leq C_p2^{-s}\|f\|_p.
\]
The middle contribution is at most
$\sum_{-B_s\leq\ell<B_s}\mathcal B_{s,\ell}f$ and has norm at most
$C_p(s+1)2^{-b_p s}\|f\|_p$ by Lemma~\ref{lem:middleband}.

For the low part the relevant $j$ form an initial interval of $[a,b]$, since
$|\mu|2^{Dj}$ increases with $j$. Write on this interval
$\Phi_{j,\mu}=\widehat K_j+\widehat{(\ee{\mu P}-1)K_j}$.
The unmodulated interval sum is a difference of two partial sums beginning
at one, and is bounded by twice the maximal function in
Lemma~\ref{lem:lifttrunc}.
On the support of $K_j$,
$|\ee{\mu P(y)}-1|\leq C_D|\mu|2^{Dj}$.
The proof of Lemma~\ref{lem:liftenvelope}, with this additional scalar factor,
therefore bounds the lifted error pointwise by
$C_D|\mu|2^{Dj}\bM f$.
The geometric series of those weights on the low interval is at most
$2^{-B_s}/(1-2^{-D})$. The entire low error is thus bounded by
$C_D2^{-B_s}\bM f$, whose $\ell^p$ norm is at most
$C_p2^{-B_s}\|f\|_p$.
For $\mu=0$ the whole interval is the unmodulated part and its error is
zero; no logarithm or band index for zero is used.

Every pointwise upper bound just obtained is independent of
$\alpha,\mu,a,b$. Take their supremum, then the $\ell^p$ norm.
Minkowski first for finite high-band sums and Fatou for their increasing
limit justify the infinite nonnegative majorant.
Finally $(s+1)2^{-b_p s}\leq C_p2^{-b_p s/2}$.
Any $0<c_p\leq\min(\theta_p,b_p/2,1)$ proves \eqref{eq:liftedSW}.
\end{proof}

\section{Kernel-side major-arc approximation and coherent grouping}\label{sec:major}
\begin{lemma}[A quantitative smoothing error]\label{lem:kernelerror}
\bp{physical\_major\_arc\_error}
Let $R=2^j$, $j\geq1$, let $1\leq L_s\leq R$, and suppose
$|\mu|R^D\leq U$. Put $h_{j,\mu}=\ee{\mu P}K_j$.
Then, on $E$ and hence on $G$,
\begin{equation}\label{eq:kernelerror}
 |h_{j,\mu}(y)-(h_{j,\mu}*_E\varphi_s)(y)|
                 \leq C(1+U)(L_s/R)B_R(y).
\end{equation}
If $\lambda=\alpha+\mu$ in $\T$, with $\alpha\in\mathcal R_s$, then the
lattice convolution kernel of
$T_{j,\lambda}-\Lift_{s,\alpha}[\Phi_{j,\mu}](D)$ is
\[
 \ee{\alpha P(y)}[h_{j,\mu}(y)-(h_{j,\mu}*_E\varphi_s)(y)].
\]
In particular the error includes its value at $y=0$.
\end{lemma}
\begin{proof}
The product and chain rules and \eqref{eq:dyadicbounds} give
\[
 \|h_{j,\mu}\|_1\leq C,\quad
 \supp h_{j,\mu}\subset\{|y|\leq2R\},\quad
 \|\nabla h_{j,\mu}\|_\infty\leq C(1+U)R^{-n-1}.
\]
Fourier inversion gives $\int\varphi_s=\chi_s(0)=1$, and its first absolute
moment is at most $CL_s$. For $|y|\leq4R$ write the difference as
$\int\varphi_s(z)[h(y)-h(y-z)]\,dz$. The mean-value bound controls it by
$C(1+U)R^{-n-1}\int |z||\varphi_s(z)|\,dz
 \leq C(1+U)L_sR^{-n-1}$.
This is \eqref{eq:kernelerror} in that region because $B_R(y)\geq cR^{-n}$.
For $|y|>4R$, $h(y)=0$ and $|y-z|\geq|y|/2$ when $z$ is in the support of
$h$. The Schwartz bound yields
\[
 |h*\varphi_s(y)|\leq C\|h\|_1 L_s^{-n}(1+|y|/L_s)^{-n-1}
                \leq C L_s|y|^{-n-1}.
\]
For $|y|>4R$, $B_R(y)$ is comparable to $R|y|^{-n-1}$, proving the remaining
region. The kernel formula follows from \eqref{eq:liftkernel} and the fact
that $P(y)\in\Z$, so equality of $\lambda$ and $\alpha+\mu$ modulo integers
is enough. At zero the original dyadic kernel vanishes, but its smoothed
counterpart need not; the difference just estimated is retained there.
\end{proof}

\begin{lemma}[The final lower scale and interval coherence]\label{lem:coherence}
\bp{major\_arc\_coherent\_intervals}
After all previous constants have been chosen, one can choose an integer
$j_0\geq j_1$ such that for every $j\geq j_0$:
\begin{enumerate}
\item $w_j<1/4$, the sequence $w_j$ is strictly decreasing, and
$w_j\leq H_j^{-2}/16$;
\item the arcs defining $X_j$ are pairwise disjoint;
\item if $1\leq s\leq\lfloor M\log_2(j+1)\rfloor$, then $v_s\leq j/2$.
\end{enumerate}
For such $s,j$ every rational center of denominator $q\in[2^{s-1},2^s)$
is allowed in $X_j$. For a fixed $s$, a fixed $\lambda\in\T$, and a finite
upper scale $J$, all $j\in[j_0,J]$ for which the unique major center at scale
$j$ belongs to $\mathcal R_s$ have the same center $\alpha$, if there are any.
Their indices form an interval beginning at
\[
 j_s=\max\bigl(j_0,\lceil2^{s/M}\rceil-1\bigr).
\]
For this interval the difference $\mu=\lambda-\alpha$ has one fixed real
representative of absolute value less than $1/4$.
\end{lemma}
\begin{proof}
The ratio of consecutive widths is
$w_{j+1}/w_j=2^{-D}(1+1/(j+1))^M$, which is less than one for sufficiently
large $j$. The limit $(j+1)^{3M}2^{-Dj}\to0$ gives
$w_j\leq (j+1)^{-2M}/16\leq H_j^{-2}/16$ and $w_j<1/4$ after increasing
the lower scale. Rational separation in Lemma~\ref{lem:divisor} then gives
disjointness, since two denominators smaller than $H_j$ have center distance
greater than $H_j^{-2}$, whereas the sum of their radii is at most
$H_j^{-2}/8$.

For an admissible $s$, \eqref{eq:vs} gives
\[
 v_s\leq1+C_*(M\log_2(j+1)+1)(1+\sqrt{j+1}).
\]
This is at most $j/2$ for all sufficiently large $j$, since
$\log(j+1)/\sqrt{j+1}\to0$. These are elementary exponential-versus-power
limits: putting $t=\log(j+1)$ reduces them to $t^ae^{-ct}\to0$, which follows,
for example, by bounding the exponential from below by a term of its power
series of integer degree greater than $a$. Choose $j_0$ above all the finite
thresholds, including $j_1$.

The inequality $H_j\geq2^s$ is equivalent to
$M\log_2(j+1)\geq s$, hence to $j\geq\lceil2^{s/M}\rceil-1$.
This proves the admissibility and the formula for $j_s$.
For any such $j$, $w_j\leq2^{-2s}/16$.
Suppose two indices in the specified set had different centers in
$\mathcal R_s$. Their torus distance would be at most the sum of the two
widths, hence at most $2^{-2s}/8$. But their denominators are smaller than
$2^s$, so rational separation gives distance greater than $2^{-2s}$.
The contradiction proves that the center $\alpha$ is fixed.
Once this center exists, it is admissible at every $j\geq j_s$.
The condition $\|\lambda-\alpha\|_\T\leq w_j$ is an initial interval of
those indices because $w_j$ decreases. Uniqueness of the major center shows
that exactly those indices occur. The restriction $j\leq J$ makes the
interval finite. Since every width is less than $1/4$, there is a unique
representative $\mu$ of the difference in $(-1/4,1/4)$, and it is independent
of $j$.
\end{proof}

\begin{proposition}[Uniform finite dyadic estimate]\label{prop:finiteSW}
\bp{finite\_dyadic\_stein\_wainger}
For every integer $J\geq j_0$ and finite-support $f$,
\begin{equation}\label{eq:finiteSW}
 \left\|\sup_{\lambda\in\T}\left|\sum_{j=j_0}^J T_{j,\lambda}f\right|\right\|_p
                  \leq C_p\|f\|_p,
\end{equation}
with a constant independent of $J$ and $f$.
\end{proposition}
\begin{proof}
At each fixed $\lambda$ split the indices into those for which
$\lambda\notin X_j$ and those for which $\lambda\in X_j$.
The maximal norm of the first sum is at most
\[
 \sum_{j=j_0}^J\left\|\sup_{\lambda\notin X_j}|T_{j,\lambda}f|\right\|_p
 \leq C_p\sum_{j\geq j_0}j^{-2}\|f\|_p
 \leq C_p\|f\|_p
\]
by \eqref{eq:minorp} and the convergent positive series.

For a major index let $\alpha=a/q$ be its unique center, let $s$ be the unique
integer with $2^{s-1}\leq q<2^s$, and take $\mu$ as in
Lemma~\ref{lem:coherence}. Because $q<H_j$ and $H_j$ is an integer power of
two, $s\leq\lfloor M\log_2(j+1)\rfloor$; no partly filled top bin occurs.
We have $|\mu|2^{Dj}\leq(j+1)^M$ and $L_s\leq2^{j/2}$.
Lemma~\ref{lem:kernelerror} therefore gives, uniformly in this choice of
$\lambda,s,\alpha,\mu$,
\[
 |T_{j,\lambda}f-\Lift_{s,\alpha}[\Phi_{j,\mu}](D)f|
       \leq C(j+1)^M2^{-j/2}(B_{2^j}*|f|).
\]
The right side has $\ell^p$ norm at most
$C_p(j+1)^M2^{-j/2}\|f\|_p$ by Young and the uniform lattice sum of $B_{2^j}$.
Its sum over $j\geq j_0$ converges, again by exponential decay. This bounds
the supremum of the total error; it does not count how many arcs are present.

For the remaining main terms, group the finite indices by $s$.
Lemma~\ref{lem:coherence} says that, for each fixed $\lambda,s$, either the
sum is empty or it is
\[
 \Lift_{s,\alpha}\left[\sum_{j=a}^b\Phi_{j,\mu}\right](D)f
\]
for one $\alpha\in\mathcal R_s$, one real $\mu$, and one finite integer
interval $1\leq a\leq b$. Its absolute value is at most $\mathcal P_s f$.
Thus the maximal main term is pointwise at most $\sum_{s\geq1}\mathcal P_s f$.
Only finitely many $s$ occur before enlarging this nonnegative sum.
Theorem~\ref{thm:liftedSW}, Minkowski, and the geometric series give
\[
 \left\|\sum_{s\geq1}\mathcal P_s f\right\|_p
 \leq C_p\sum_{s\geq1}2^{-c_p s}\|f\|_p
 \leq C_p\|f\|_p.
\]
Taking increasing finite partial sums and applying Fatou justifies the
infinite majorant. Adding minor terms, smoothing errors and major terms
proves \eqref{eq:finiteSW}.
\end{proof}

\section{The operator on all of \texorpdfstring{$\ell^p$}{lp} and the final theorem}\label{sec:final}
\begin{lemma}[Absolute convergence and the scalar supremum]\label{lem:absolute}
\bp{absolute\_convergence\_and\_parameter\_continuity}
For every $1<p<\infty$ the restricted kernel
$K\ind_{G\setminus\{0\}}$ belongs to $\ell^{p'}(G)$.
For $f\in\ell^p(G)$ and each $x\in G$ the series
\[
 T_\lambda f(x)=\sum_{y\in G\setminus\{0\}}
                    \ee{\lambda P(y)}K(y)f(x-y)
\]
is absolutely convergent, uniformly in $\lambda\in\R$.
It is continuous and one-periodic in $\lambda$; its absolute value attains a
maximum on $\T$, equal to the supremum over rational parameters.
For any two $\ell^p$ inputs,
\begin{equation}\label{eq:pointwisecontinuity}
 \sup_\lambda|T_\lambda(f-g)(x)|
            \leq\|K\|_{\ell^{p'}(G\setminus\{0\})}\|f-g\|_p.
\end{equation}
\end{lemma}
\begin{proof}
Compactness of the sphere bounds $|\Omega|$, hence $|K(y)|\leq C|y|^{-n}$.
There are at most $C_n2^{kn}$ lattice points with $2^k\leq|y|<2^{k+1}$.
Their contribution to the $p'$th power of the kernel norm is at most
$C2^{-kn(p'-1)}$. Its sum for $k\geq0$ converges because $p'>1$; this covers
all nonzero lattice points since their norm is at least one.
H\"older now gives
$\sum_{y\ne0}|K(y)f(x-y)|\leq\|K\|_{p'}\|f\|_p$.
The tails of this positive convergent series bound the tails uniformly in
$\lambda$. Finite sums are continuous, so their uniform limit is continuous:
for a specified error first make both tails smaller than a third of it,
then use continuity of the finite sum. Since $P(y)\in\Z$, each summand has
period one. Compactness of $[0,1]$ and the extreme-value theorem give the
maximum. Density of rational numbers and continuity show that its value is
the supremum over rational parameters. Applying H\"older to $f-g$ proves
\eqref{eq:pointwisecontinuity}.
\end{proof}

\begin{theorem}[The combined discrete Stein--Wainger bound]\label{thm:main}
\bp{discrete\_stein\_wainger\_lp}
Let $n,d\geq1$ be integers. Let $\Omega:E\setminus\{0\}\to\C$ be smooth,
homogeneous of degree zero, and have integral zero on $S^{n-1}$.
Set $K(y)=\Omega(y)|y|^{-n}$ for $y\ne0$ and
$P(y)=(y_1^2+\cdots+y_n^2)^d$.
For every real $p$ with $1<p<\infty$ there is a finite constant
$C_{p,n,d,\Omega}$ such that, for all $f\in\ell^p(\Z^n)$,
\begin{equation}\label{eq:MAIN}
 \left\|\sup_{\lambda\in\R}\left|
       \sum_{y\in\Z^n\setminus\{0\}}
               \ee{\lambda P(y)}K(y)f(x-y)\right|\right\|_{\ell^p_x}
       \leq C_{p,n,d,\Omega}\|f\|_{\ell^p}.
\end{equation}
The sum is the absolutely convergent sum in Lemma~\ref{lem:absolute}.
The supremum is over a scalar parameter, and only finite $p>1$ are asserted.
\end{theorem}
\begin{proof}
All choices and estimates above were made for an arbitrary fixed
$p\in(1,\infty)$. For finite-support $f$, the telescoping identity gives at
nonzero lattice points
$K(y)=\eta(y)K(y)+\sum_{j\geq1}K_j(y)$.
The initial kernel
\[
 k_{\mathrm{init}}(y)=|\eta(y)K(y)|\ind_{y\ne0}
                           +\sum_{1\leq j<j_0}|K_j(y)|
\]
has finite $\ell^1$ norm: the first term is supported on finitely many lattice
points, and each of the finitely many remaining terms has uniformly bounded
$\ell^1$ norm. Its positive convolution bounds the maximal contribution of
these initial terms, by Young.
For every $J\geq j_0$, Proposition~\ref{prop:finiteSW} bounds the remaining
finite dyadic sum with a constant independent of $J$.

At each fixed $x$, only finitely many $y$ have $f(x-y)\ne0$. For each such
nonzero $y$, only finitely many $j$ have $K_j(y)\ne0$, by the annular support.
Therefore the initial contribution plus the finite dyadic sums equals
$T_\lambda f(x)$ for all sufficiently large $J$, simultaneously for all
$\lambda$. Its supremum is eventually exactly the desired supremum at that
$x$. Fatou transfers the uniform finite-sum estimate to \eqref{eq:MAIN} for
finite-support $f$.

For an arbitrary $f\in\ell^p$, take its restrictions $f_R$ to increasing
finite cubes. Lemma~\ref{lem:finite} gives $\|f_R-f\|_p\to0$ and
$\|f_R\|_p\to\|f\|_p$. By \eqref{eq:pointwisecontinuity},
$\sup_\lambda|T_\lambda(f_R-f)(x)|\to0$ at every $x$.
The inequality
$|\sup_\lambda|T_\lambda f_R(x)|-\sup_\lambda|T_\lambda f(x)||
 \leq\sup_\lambda|T_\lambda(f_R-f)(x)|$
therefore gives convergence of the maximal functions pointwise.
Fatou and the already established bounds for $f_R$ yield
\[
 \|\sup_\lambda|T_\lambda f|\|_p
 \leq\liminf_R C_{p,n,d,\Omega}\|f_R\|_p
 =C_{p,n,d,\Omega}\|f\|_p.
\]
This is the asserted theorem. No endpoint at $p=1$ or $p=\infty$ is obtained
by this argument or included in the statement.
\end{proof}

\section{Implementation contracts and the dependency manifest}\label{sec:implementation}
This section restates the already proved mathematics as implementation
contracts. It does not introduce additional analytic hypotheses. The local
identifiers printed with the lemmas are proposed declaration names; none is
a claim that an identically named declaration already exists in mathlib.

\subsection{Concrete types and representations}
Use \ml{Fin n} for coordinates, \ml{Fin n -> Int} for the lattice, and
\ml{EuclideanSpace Real (Fin n)} for real space. Here \ml{->} denotes the function arrow in a type specification,
not a claim that this printed text is compiled source. The latter
choice is important: the default norm on a function type is not the
Euclidean norm used by the phase. Define the lattice embedding coordinatewise
and evaluate its Euclidean norm only after that embedding. Counting measure
on the lattice has no nonempty null sets, so equality almost everywhere is
pointwise equality there.

Define an integer-valued polynomial first,
\[
 P_\Z(y)=\left(\sum_{i\in\mathrm{Fin}\ n} y_i^2\right)^d\in\Z,
\]
and a real-valued polynomial by the same expression on Euclidean space.
The identity $P_\R(\iota y)=\operatorname{cast}(P_\Z(y))$ follows by moving
the ring homomorphism through the finite sum, square and final power.
The identity $P_\R(x)=\|x\|^{2d}$ follows from
$\|x\|^2=\sum_i x_i^2$ and $(\|x\|^2)^d=\|x\|^{2d}$.
This avoids reasoning about an arbitrary real power of a lattice norm in
all arithmetic lemmas. Lemma~\ref{lem:radialcoeff} supplies the coefficient
and congruence interfaces.

Represent $\Omega$ as a total function on Euclidean space with an irrelevant
chosen value at zero, together with smoothness on the complement of zero,
positive-dilation invariance there, and its zero sphere integral. Define
$K(0)=0$ by a case split. None of the differentiability statements about $K$
includes zero. In contrast $K_j$ is globally smooth because it vanishes on a
neighborhood of zero. The sphere measure can be implemented using the cone
construction of Lemma~\ref{lem:dyadic}; that lemma proves its equivalence with
the surface measure in the target specification.

Represent a rational center by positive $q$, $0\leq a<q$, and $\gcd(a,q)=1$,
then map $a/q$ to \ml{AddCircle (1 : Real)}. In a finite rational set use the
image finite set, rather than a list that may count duplicate torus points.
The full grids $T[q]$ can be represented by \ml{Fin q} in each coordinate.
The block decomposition is the explicit finite bijection in
Definition~\ref{def:blocks}, not a heuristic separation of prime factors.
For $q=1$ it has one element and all finite character identities remain valid.

\subsection{Norms, finite maxima, and operator interfaces}
For a real exponent $p>1$ the library exponent is its image in
$\mathbb R_{\geq0}\cup\{\infty\}$, implemented by \ml{ENNReal.ofReal p}.
Prove first that this exponent is at least one and is not infinity.
Use \ml{MeasureTheory.MemLp} and the extended norm while establishing
membership. Only then pass to \ml{MeasureTheory.Lp} and its ordinary norm.
In particular, never use conversion of an infinite extended norm to a real
number as a substitute for a finiteness proof.

A safe finite version of every lattice maximal estimate is
\[
 \sum_{x\in E_0}\left(\max_{a\in F}|T_a f(x)|\right)^p
             \leq C^p\sum_{x\in G}|f(x)|^p,
\]
for arbitrary finite output sets $E_0$ and finite parameter sets $F$.
Include zero in the maximum to give the empty set value zero.
Lemma~\ref{lem:finite} passes from these bounds to the full supremum.
For the final operator, Lemma~\ref{lem:absolute} already supplies pointwise
boundedness in the parameter; its real supremum is therefore a legitimate
conditionally complete-lattice supremum. It also identifies that supremum
with the countable rational one.

Implement the Euclidean and torus multipliers first on $L^2$ by the Fourier
isometries. An assertion that a bounded measurable symbol is an $L^p$
multiplier of norm at most $A$ means that its $L^2$ action on Schwartz inputs
satisfies the stated $L^p$ inequality. The density and extension lemma then
constructs its bounded $L^p$ action. When two such extensions agree on
Schwartz functions, their difference vanishes on a dense set and hence on
the entire space by continuity. On the lattice use finite-support functions
as the dense common domain. This is the identification used for all kernel,
periodization and multiplier factorizations in the proof.

Hilbert-valued estimates only require finite products of $\C$. For an
index set $I$ carry the Euclidean norm
$(\sum_{i\in I}|z_i|^2)^{1/2}$ explicitly, rather than the default supremum
norm on a function type. Lemmas~\ref{lem:interpolation} and
\ref{lem:khintchine} apply uniformly to these finite products; no theorem
about infinite-dimensional vector-valued singular integrals is needed.

For a finite family, implement linearization by an ordered argmax, with ties
resolved by the least index. The selected map depends on the input used to
choose it, but after selection it is held fixed while forming the linear
operator and its adjoint. The uniform $TT^*$ estimates quantify over every
such fixed selection. Thus there is no claim that the maximal operator
itself is linear. In Euclidean space the selection is measurable because its
fibers are finite intersections of comparisons of measurable absolute values.
On the lattice measurability is automatic.

\subsection{The final declaration contract}
The mathematical contract for \ml{discrete_stein_wainger_lp} is the following;
it is a specification, not purportedly compiled Lean syntax.
\begin{quote}
For positive natural numbers $n,d$, admissible $\Omega$, and real $p>1$,
there exists a real $C\geq0$ such that for every function
$f:(\mathrm{Fin}\ n\to\Z)\to\C$ with $\operatorname{MemLp}(f,p,\mathrm{count})$,
the maximal function $\mathcal C f$ belongs to
$L^p(\mathrm{count};\R)$ and
\[
 \operatorname{eLpNorm}(\mathcal C f,p,\mathrm{count})
 \leq \operatorname{ofReal}(C)\,
                   \operatorname{eLpNorm}(f,p,\mathrm{count}).
\]
Here the defining series is summable by Lemma~\ref{lem:absolute}, the phase
uses the cast of $P_\Z$, and $\mathcal C f$ is the scalar-parameter supremum
of its complex absolute value.
\end{quote}
The membership conclusion is essential: it excludes a vacuous inequality
obtained by coercing an infinite norm to zero. Since counting measure has
only the empty null set, the resulting representative agrees at every
lattice point with the series-defined maximal function.

\subsection{Constant choice and termination of dependencies}
The constants are selected in this order:
\[
 (n,d,\Omega,p)\ ;\ r\ ;\ \kappa\ ;\ (u,v,b,M,j_1)\ ;\ \rho\ ;\
 C_*\ ;\ (\gamma,\delta_p)\ ;\ C_1\ ;\ (\theta_p,b_p,\sigma_p,C_2,c_p)\ ;\ j_0.
\]
The exponent $\sigma_p$ can already be read from
Lemma~\ref{lem:continuousband}; listing it at its use does not make it depend
on later constants. All parameter choices are justified by explicit
inequalities in the corresponding constructions. $M$ is chosen only after
the minor-arc counting exponents, whereas $j_0$ is chosen only after the
smoothing scales. This avoids the circular condition of choosing a cutoff
narrowness from a lower-scale threshold that itself depends on that cutoff.
The positive gains $\gamma,\delta_p,\theta_p,b_p,c_p$ need not be optimized;
the final summation uses only their strict positivity.

The inverse-sum proof uses a finite induction on the degree, not a recursive
appeal to the final theorem. Its polynomial-loss convention can be expanded
mechanically: associate a pair $(C,A)$ to a bound $C\delta^{-A}$, add
exponents when multiplying bounds, take the maximum exponent and add
constants for sums, and multiply both exponents and logarithms of constants
for fixed powers. Each occurrence of ``sufficiently large $N$'' in that
finite induction is an explicit comparison $N\geq C\delta^{-A}$.
The complementary small-$N$ case is handled at the start of the induction
by the trivial torus-distance bound $1/2$. No limiting induction or
ineffective compactness theorem is hidden in that convention.

\subsection{Integrability and limit register}
There are four distinct convergence operations, each with an already proved
majorant or norm bound. Finite lattice annular kernels and finite rational
sums need no convergence argument. Infinite sums involving the Schwartz
cutoffs have absolute polynomial-decay majorants, as in
Lemmas~\ref{lem:periodic} and \ref{lem:liftkernel}. The whole real singular
kernel and the signed smoothing errors use the $L^2$/Fatou limit in
Lemma~\ref{lem:signedK} and the absolutely $L^1$-convergent series
\eqref{eq:EJseries}. Maxima over infinitely many scales or parameter values
use finite-family estimates and Lemma~\ref{lem:finite}, with countable
parameter density only on Euclidean space. Finally the original operator
on all of $\ell^p$ uses the absolute pointwise series and the uniform
parameter estimate \eqref{eq:pointwisecontinuity}.

The integral manipulations in the long-truncation averaging argument are
finite averages of nonnegative $p$th powers followed by lattice sums.
They are justified by Tonelli even before their bound is known. The changes
of variables are the explicit finite-residue bijections in that proof.
They are not exchanges of a supremum with an oscillatory infinite integral.

\subsection{Dependency index}
The following index records the principal direct numbered dependencies.
The elementary library boundary F1--F5, preceding definitions, and elementary
algebraic manipulations are implicit. References always point to earlier
proof nodes; the table itself is only a navigation aid.

\begingroup
\small
\setlength{\tabcolsep}{5pt}
\renewcommand{\arraystretch}{1.2}
\begin{longtable}{@{}p{.48\textwidth}p{.47\textwidth}@{}}
\toprule
Proof node & Principal earlier dependencies \\
\midrule
\endfirsthead
\toprule
Proof node & Principal earlier dependencies \\
\midrule
\endhead
\midrule
\multicolumn{2}{r}{\footnotesize Continued on the next page.}\\
\endfoot
\bottomrule
\endlastfoot
\ref{lem:finite}. Finite reductions & F3, F4 \\
\ref{lem:young}. Convolution, duality, subsequences & F3, F4 \\
\ref{lem:interpolation}. Distributional interpolation & \ref{lem:young} \\
\ref{lem:maximal}. Maximal function and envelopes & \ref{lem:interpolation} \\
\ref{lem:ttstar}. Hilbert-space kernel bounds & \ref{lem:finite}, \ref{lem:young} \\
\ref{lem:khintchine}. Randomization and square functions & F1, F3, F4 \\
\ref{lem:binary}. Binary decomposition & F1 \\
\ref{lem:radialcoeff}. Radial polynomial coefficients & F1, F2 \\
\ref{lem:divisor}. Divisor and progression counts & F1 \\
\ref{lem:multiples}. Approximate multiples & F1 \\
\ref{lem:inverseweyl}. Inverse exponential sum & \ref{lem:multiples} \\
\ref{lem:weylcor}. Weighted slicing and rational gain & \ref{lem:inverseweyl} \\
\ref{lem:gauss}. Quadratic sums & F1 \\
\ref{lem:periodic}. Periodic averaging & \ref{lem:divisor} \\
\ref{lem:dyadic}. Polar integration and dyadic kernels & F1--F4 \\
\ref{lem:cz}. Integrable singular kernels & \ref{lem:young}, \ref{lem:interpolation} \\
\ref{lem:signedK}. Signed dyadic multipliers & \ref{lem:dyadic}, \ref{lem:cz}, \ref{lem:young} \\
\ref{con:phi}. Positive low-frequency averaging & F4, F5 \\
\ref{lem:smootherrors}. Signed smoothing errors & \ref{lem:dyadic}, \ref{lem:cz}, \ref{lem:signedK}, \ref{con:phi} \\
\ref{lem:oscint}. Polynomial oscillatory integrals & F1, F2 \\
\ref{lem:radialFourier}. Radial stationary bounds & \ref{lem:radialcoeff}, \ref{lem:dyadic} \\
\ref{lem:continuousband}. Continuous annular decay & \ref{lem:radialcoeff}, \ref{lem:dyadic}, \ref{lem:oscint}, \ref{lem:ttstar}, \ref{lem:maximal}, \ref{lem:interpolation} \\
\ref{lem:fourierbook}. Fourier bookkeeping & F1, F3--F5 \\
\ref{lem:sampling}. Sampling and full grids & \ref{lem:young}, \ref{lem:fourierbook} \\
\ref{def:blocks}. Coprime-block rational sets & \ref{lem:divisor} \\
\ref{lem:superorth}. Singleton superorthogonality & \ref{lem:khintchine} \\
\ref{lem:sunflower}. Weighted sunflower inequality & F1, F3 \\
\ref{lem:blockorthog}. Block and numerator orthogonality & \ref{def:blocks}, \ref{lem:fourierbook} \\
\ref{lem:normdescription}. Rational norm description & \ref{lem:superorth}, \ref{lem:sunflower}, \ref{lem:blockorthog} \\
\ref{lem:companion}. Companion projection bound & \ref{def:blocks}, \ref{lem:blockorthog}, \ref{lem:interpolation}, \ref{lem:fourierbook} \\
\ref{thm:IW}. Rational-frequency transference & \ref{lem:sampling}, \ref{lem:normdescription}, \ref{lem:companion}, \ref{lem:khintchine}, \ref{lem:young}, \ref{lem:interpolation} \\
\ref{con:goodset}. Good denominator set & \ref{def:blocks}, \ref{thm:IW} \\
\ref{lem:minorL2}. Minor-arc Hilbert-space decay & \ref{lem:dyadic}, \ref{lem:radialcoeff}, \ref{lem:weylcor}, \ref{lem:multiples}, \ref{lem:ttstar}, \ref{lem:finite} \\
\ref{con:parameters}. Exponent and logarithmic scale & \ref{lem:minorL2}, \ref{lem:interpolation}, \ref{lem:dyadic}, \ref{con:goodset} \\
\ref{lem:liftkernel}. Exact physical lifting identity & \ref{lem:fourierbook}, \ref{lem:radialcoeff} \\
\ref{lem:arithmax}. Arithmetic maximal decay & \ref{lem:liftkernel}, \ref{lem:periodic}, \ref{lem:weylcor}, \ref{lem:gauss}, \ref{lem:divisor}, \ref{lem:maximal}, \ref{lem:ttstar}, \ref{lem:interpolation} \\
\ref{lem:liftsquare}. Lifted square estimates & \ref{lem:arithmax}, \ref{thm:IW}, \ref{con:goodset}, \ref{lem:khintchine} \\
\ref{lem:liftenvelope}. Lifted positive envelopes & \ref{lem:liftkernel}, \ref{lem:maximal}, \ref{lem:dyadic} \\
\ref{lem:lifttrunc}. Lifted maximal truncation & \ref{lem:binary}, \ref{lem:signedK}, \ref{lem:smootherrors}, \ref{con:phi}, \ref{lem:liftsquare}, \ref{lem:sampling}, \ref{lem:maximal}, \ref{lem:fourierbook} \\
\ref{lem:middleband}. Uniform band gain & \ref{lem:liftsquare}, \ref{lem:radialFourier}, \ref{lem:liftenvelope}, \ref{lem:interpolation} \\
\ref{lem:highband}. High-band analytic gain & \ref{lem:continuousband}, \ref{lem:sampling} \\
\ref{thm:liftedSW}. Lifted partial sums & \ref{lem:lifttrunc}, \ref{lem:middleband}, \ref{lem:highband}, \ref{lem:liftenvelope}, \ref{lem:maximal}, \ref{lem:finite} \\
\ref{lem:kernelerror}. Kernel-side approximation & \ref{lem:dyadic}, \ref{lem:liftkernel} \\
\ref{lem:coherence}. Final scale and coherent arcs & \ref{lem:divisor}, \ref{con:parameters} \\
\ref{prop:finiteSW}. Uniform finite dyadic estimate & \ref{con:parameters}, \ref{lem:kernelerror}, \ref{lem:coherence}, \ref{thm:liftedSW}, \ref{lem:young}, \ref{lem:finite} \\
\ref{lem:absolute}. Pointwise absolute convergence & \ref{lem:dyadic}, \ref{lem:young} \\
\ref{thm:main}. The final maximal theorem & \ref{prop:finiteSW}, \ref{lem:absolute}, \ref{lem:dyadic}, \ref{lem:finite}, \ref{lem:young} \\
\end{longtable}
\endgroup

\subsection{Specification audit and library entry points}
The two source papers are used to specify the target operator and its
hypotheses, not as unproved lemmas in the dependency graph. In particular,
the $\ell^2$ case is derived above as part of the finite-$p$ argument; it is
not imported from the first paper. The real annular maximal estimate,
arithmetic inverse estimate, coprime-block transference, and lifted maximal
truncation each have their own proof before use.

Three features of the specification are worth preserving in an implementation.
First, the parameter is a scalar real number; its reduction modulo one uses
that $P_\Z(y)$ is an integer. Second, homogeneity and zero spherical mean
are used for \emph{every} radially modulated annulus in
\eqref{eq:radialcancel}, not merely for the unmodulated kernel. Third, the
conclusion is for finite $p>1$. The argument does not assert an $\ell^\infty$
endpoint. In particular, one must not replace maximal-function control of
annular envelopes by an assertion that the supremum of all annular kernels
is summable.

\textbf{Target specifications.}
Ben Krause and Joris Roos,
\emph{Discrete analogues of maximally modulated singular integrals of
Stein--Wainger type}, version 3, 17 March 2021:
\url{https://arxiv.org/abs/1907.00405v3}.
The same authors,
\emph{Discrete analogues of maximally modulated singular integrals of
Stein--Wainger type: $\ell^p$ bounds for $p>1$}, version 2, 23 July 2023:
\url{https://arxiv.org/abs/2107.14616v2}.
These bibliographic entries identify the problem, and are not additional
proof dependencies.

\textbf{Mathlib boundary documentation.}
The following entry points locate the imported Fourier, approximation, and
integral-norm material. Exact declarations named in Section~1 were checked
against these pages. Their hypotheses still have to be instantiated with
the measures, normalizations, and finite-dimensional spaces of this document.
\begin{flushleft}\small
\url{https://leanprover-community.github.io/mathlib4_docs/Mathlib/Analysis/Fourier/LpSpace.html}\par
\url{https://leanprover-community.github.io/mathlib4_docs/Mathlib/Analysis/Fourier/AddCircleMulti.html}\par
\url{https://leanprover-community.github.io/mathlib4_docs/Mathlib/Analysis/Distribution/SchwartzSpace/Basic.html}\par
\url{https://leanprover-community.github.io/mathlib4_docs/Mathlib/MeasureTheory/Integral/MeanInequalities.html}
\end{flushleft}
The elementary calculus, integration, algebra, and bump-function facts listed
in F1--F4 are the rest of the permitted library boundary. No use of a
literature theorem is implicit in a local blueprint identifier.

\textbf{Verification boundary.}
This is a proposed mathematical formalization blueprint. The dependency order,
finite-family reductions, convergence obligations, and declaration contracts
are explicit. It is not a Lean proof certificate, and no claim is made that
Lean has checked the estimates or that the proposed declaration names exist
in mathlib. A formal implementation should pin a mathlib revision, discharge
the nodes in the displayed order, and only then export the final declaration.

\end{document}
